\documentclass[11pt]{article}

\usepackage[margin=1in]{geometry}
\usepackage{amsmath,amssymb,amsthm,mathtools}
\usepackage{booktabs}
\usepackage{enumitem}
\usepackage{hyperref}
\usepackage{microtype}
\usepackage{xcolor}

\hypersetup{colorlinks=true,linkcolor=blue,citecolor=blue,urlcolor=blue}
\allowdisplaybreaks

\newtheorem{proposition}{Proposition}
\newtheorem{lemma}{Lemma}
\newtheorem{definition}{Definition}

\newcommand{\R}{\mathbb R}
\newcommand{\E}{\mathbb E}
\newcommand{\dd}{\,\mathrm d}
\newcommand{\given}{\,\middle|\,}
\newcommand{\ind}{\mathbf 1}
\newcommand{\bfeta}{\boldsymbol\eta}
\newcommand{\bfalpha}{\boldsymbol\alpha}
\newcommand{\bfxi}{\boldsymbol\xi}
\newcommand{\bfz}{\boldsymbol z}
\newcommand{\bfu}{\boldsymbol u}
\newcommand{\bfr}{\boldsymbol r}
\newcommand{\diag}{\operatorname{diag}}
\newcommand{\vecop}{\operatorname{vec}}
\newcommand{\tr}{\operatorname{tr}}

\newcommand{\widehatell}{\widehat\ell}
\newcommand{\widehatZ}{\widehat Z}
\newcommand{\widehatq}{\widehat q}
\newcommand{\widehatp}{\widehat p}
\newcommand{\dotwidehatZ}{\dot{\widehat Z}}
\newcommand{\dotwidehatq}{\dot{\widehat q}}
\newcommand{\dotwidehatp}{\dot{\widehat p}}
\newcommand{\cond}{\,|\,}
\newcommand{\shape}{\operatorname{shape}}

\title{An Integrated Readable Companion to Zhao--Cui Tensor-Train Sequential Filtering and Fixed-Branch Gradients}
\author{BayesFilter P22 Technical Note}
\date{2026-06-02}

\begin{document}
\maketitle

\begin{abstract}
This P22 note is an integrated readable companion built from the P20
source-order reconstruction and the P21 implementation-orientation supplement.
It does not summarize, shorten, or replace P20.  Instead, it preserves P20's
Zhao--Cui annotation and fixed-branch derivations, then inserts additional
orientation, shape contracts, carried-filter storage details, and
finite-difference reporting rules at the points where a reader or later
implementer needs them.  The document is intentionally generous: it favors
explicit equations, dimensions, and implementation meaning over compactness.
\end{abstract}

\tableofcontents

\section{Reader Contract, Non-Summary Rule, And Notation}

This note uses ``Zhao--Cui anchor'' sentences to identify where a displayed
formula comes from, but the anchor is not the explanation.  Every important
formula is restated in BayesFilter notation, derived from the previous
formula, interpreted for filtering or transport, and tied to an implementation
object.

The construction rule for P22 is
\[
    \text{P22}
    =
    \text{P20 mathematical spine}
    +
    \text{P21 reader orientation and implementation specification},
    \tag{P22-C1}
\]
not
\[
    \text{P22}
    =
    \text{a shorter summary of P20}.
    \tag{P22-C2}
\]
The source-order Zhao--Cui annotation remains the spine through Sections
1--3 and 5.  The added material has a narrower purpose: whenever an equation
is difficult to read, it states what is being stored, what is being integrated,
what is being differentiated, which branch choices are fixed, and what would
break the mathematical claim.

Zhao--Cui write the unknown parameter as \(\theta\).  The first part of this
note follows that convention when reconstructing their paper.  The
BayesFilter derivative extension later writes the external differentiated
parameter as \(\beta\), because there \(\theta\) would otherwise mean both a
random coordinate and the argument of the likelihood function.

The state dimension is \(m\), observation dimension is \(n\), and parameter
dimension is \(d\).  Thus
\[
    x_t=(x_{t,1},\ldots,x_{t,m})\in\mathcal X\subseteq\R^m,
    \qquad
    y_t=(y_{t,1},\ldots,y_{t,n})\in\mathcal Y\subseteq\R^n,
\]
and
\[
    \theta=(\theta_1,\ldots,\theta_d)\in\Theta\subseteq\R^d.
\]
For coordinate groups, the paper uses
\[
    x_{t,<j}=(x_{t,1},\ldots,x_{t,j-1}),\qquad
    x_{t,>j}=(x_{t,j+1},\ldots,x_{t,m}),
    \tag{N1}
\]
and similarly
\[
    x_{t,\le j}=(x_{t,1},\ldots,x_{t,j}),\qquad
    x_{t,\ge j}=(x_{t,j},\ldots,x_{t,m}).
    \tag{N2}
\]
The point of this notation is conditional density bookkeeping.  A lower
triangular map conditions coordinate \(j\) on \(x_{<j}\); an upper triangular
map conditions it on \(x_{>j}\).

All densities in the reconstruction are assumed absolutely continuous with
respect to the relevant Lebesgue measure.  Zhao--Cui use \(p\) for normalized
posterior densities and \(\pi\) for unnormalized densities.  Their
approximations are \(\widehat p\) and \(\widehat\pi\).  When \(p\) and
\(\widehat p\) are normalized densities, the Hellinger distance is
\[
    D_H(p,\widehat p)
    =
    \left[
    \frac12
    \int
    \bigl(\sqrt{p(r)}-\sqrt{\widehat p(r)}\bigr)^2\dd r
    \right]^{1/2}.
    \tag{N3}
\]
This distance appears because the squared-TT method approximates square roots
of densities.

If \(S:\mathcal X\to\mathcal U\) is a diffeomorphism and \(X\sim p\), then
the density of \(U=S(X)\) is the pushforward
\[
    S_\#p(u)
    =
    p(S^{-1}(u))\left|\det\nabla S^{-1}(u)\right|.
    \tag{N4}
\]
If \(U\sim\eta\), then the density of \(X=S^{-1}(U)\) is the pullback
\[
    S^\#\eta(x)
    =
    \eta(S(x))\left|\det\nabla S(x)\right|.
    \tag{N5}
\]
These two identities are the entire mathematical engine behind the
Knothe--Rosenblatt maps and the preconditioning maps in Section 5.  The
implementation stores an evaluator for \(S\), an evaluator for \(S^{-1}\), and
the relevant log-Jacobian determinant when the map is used as a density
transform.

\section{Reader Orientation: Five Mathematical Objects}

The Zhao--Cui method is easier to follow if the notation is organized around
five objects.  They are not five separate algorithms; they are five views of
one sequential density approximation:
\[
    q_t,\qquad \phi_t,\qquad C_{t,k},\qquad \widehat Z_t,
    \qquad \partial_\beta\log\widehat Z_t .
    \tag{P22-O1}
\]
The first three describe the approximated density.  The fourth is the scalar
contributed to the likelihood.  The fifth is the sensitivity needed by the
BayesFilter fixed-branch extension.

\subsection*{Object 1: the unnormalized density}

Bayes' rule begins with an unnormalized numerator:
\[
    \text{posterior density}
    =
    \frac{\text{likelihood}\times\text{prior}}
    {\int \text{likelihood}\times\text{prior}}.
    \tag{P22-O2}
\]
In filtering, the prior at time \(t\) is the previous filter propagated through
the transition model.  A typical adjacent-state target has the form
\[
    q_t(x_t,x_{t-1};\beta)
    =
    \widehat p_{t-1}(x_{t-1};\beta)
    f_\beta(x_t\mid x_{t-1})
    g_\beta(y_t\mid x_t).
    \tag{P22-O3}
\]
The implementation question is more concrete than the notation suggests:
which coordinates are used, and at which fitting points is \(q_t\) evaluated?
A fixed branch stores
\[
    Z_{\rm fit}
    =
    \begin{bmatrix}
    z_1^{(1)} & z_2^{(1)}\\
    \vdots & \vdots\\
    z_1^{(N)} & z_2^{(N)}
    \end{bmatrix},
    \qquad
    \shape(Z_{\rm fit})=(N,2),
    \tag{P22-O4}
\]
and a target vector
\[
    \widetilde q_t^{\rm fit}
    =
    \left(
    \widetilde q_t(z^{(1)};\beta),\ldots,
    \widetilde q_t(z^{(N)};\beta)
    \right)^\top,
    \qquad
    \shape(\widetilde q_t^{\rm fit})=(N,).
    \tag{P22-O5}
\]

\subsection*{Object 2: the square-root approximation}

A direct TT fit of \(q_t\) can be negative.  Zhao--Cui therefore fit a square
root and square it back:
\[
    \phi_t(z;\beta)\approx e^{c_t/2}\sqrt{\widetilde q_t(z;\beta)},
    \tag{P22-O6}
\]
\[
    \widehat q_t(z;\beta)
    =
    e^{-c_t}\phi_t(z;\beta)^2+\tau_t\lambda_t(z).
    \tag{P22-O7}
\]
The fitted values are
\[
    y_t^{\rm fit}[j]
    =
    e^{c_t/2}\sqrt{\widetilde q_t(z^{(j)};\beta)},
    \qquad
    \shape(y_t^{\rm fit})=(N,),
    \tag{P22-O8}
\]
and the fixed-branch derivative values are
\[
    \dot y_t^{\rm fit}[j]
    =
    e^{c_t/2}
    \frac{\dot{\widetilde q}_t(z^{(j)};\beta)}
    {2\sqrt{\widetilde q_t(z^{(j)};\beta)}}.
    \tag{P22-O9}
\]

\subsection*{Object 3: the tensor-train cores}

For two coordinates, the fitted square root can be written
\[
    \phi_t(z_1,z_2;\beta)
    =
    \sum_{r=1}^{R}
    \left[
    \sum_{\ell=0}^{p-1}C_1[0,\ell,r]b_1(z_1)[\ell]
    \right]
    \left[
    \sum_{m=0}^{p-1}C_2[r,m,0]b_2(z_2)[m]
    \right].
    \tag{P22-O10}
\]
The stored arrays are
\[
    \shape(C_1)=\shape(\dot C_1)=(1,p,R),
    \qquad
    \shape(C_2)=\shape(\dot C_2)=(R,p,1).
    \tag{P22-O11}
\]
The branch freezes the fitting rule, not the fitted numerical core values:
\[
    C_k(\beta;B)
    =
    \text{core returned by the fixed fitting equations at }\beta.
    \tag{P22-O12}
\]

\subsection*{Object 4: the normalizer}

The evidence increment is the mass of the declared approximate density:
\[
    \widehat Z_t(\beta)
    =
    e^{-c_t}R_t(\beta)+\tau_t,
    \qquad
    R_t(\beta)=\int\phi_t(z;\beta)^2\dd z.
    \tag{P22-O13}
\]
In TT form, this mass is computed by contractions rather than by a full tensor
grid.  For the two-coordinate case,
\[
    S_2[r,s]
    =
    \sum_{\ell,m}
    C_2[r,\ell,0]B_2[\ell,m]C_2[s,m,0],
    \tag{P22-O14}
\]
\[
    R_t
    =
    \sum_{r,s,\ell,m}
    C_1[0,\ell,r]B_1[\ell,m]C_1[0,m,s]S_2[r,s].
    \tag{P22-O15}
\]

\subsection*{Object 5: the derivative of the log normalizer}

The derivative used by the fixed-branch likelihood is
\[
    \partial_\beta\log\widehat Z_t
    =
    \frac{\dot{\widehat Z}_t}{\widehat Z_t}.
    \tag{P22-O16}
\]
Since \(c_t,\tau_t,\lambda_t\) are frozen in the fixed branch,
\[
    \dot{\widehat Z}_t=e^{-c_t}\dot R_t.
    \tag{P22-O17}
\]
The derivative of the carried filter is equally important because it becomes
part of the next target:
\[
    \dot{\widehat p}_t(z_t;\beta)
    =
    \frac{\dot a_t(z_t;\beta)}{\widehat Z_t(\beta)}
    -
    \frac{a_t(z_t;\beta)\dot{\widehat Z}_t(\beta)}
    {\widehat Z_t(\beta)^2}.
    \tag{P22-O18}
\]
This front orientation is not a replacement for the derivation below.  It is
the small map used to read the longer source-order reconstruction.

\section{What Problem Is Being Solved?}

\paragraph{Source unit S1-overview.}
Zhao--Cui begin with a hidden Markov state, observations, and possibly unknown
static parameters.  The question answered here is: what density must a
sequential algorithm update when new data arrive?

The data arrive in time order.  At time \(t\) there is a hidden state
\[
    x_t\in \R^{m_x}
\]
and an observation
\[
    y_t\in \R^{m_y}.
\]
There may also be a static unknown parameter
\[
    \alpha\in \R^{m_\alpha}.
\]
The filtering problem is to update the conditional density of the current state
and parameter when \(y_t\) arrives:
\[
    p(x_t,\alpha\mid y_{1:t}).
\]
The path problem is to update the density of the whole trajectory and parameter:
\[
    p(x_{0:t},\alpha\mid y_{1:t}).
\]
The smoothing problem asks for an earlier state after later data have been seen:
\[
    p(x_k\mid y_{1:t}),\qquad k<t.
\]

The bottleneck is always an integral.  At time \(t-1\) suppose we know, or have
approximated,
\[
    p(x_{t-1},\alpha\mid y_{1:t-1}).
\]
After seeing \(y_t\), Bayes' rule introduces the unnormalized three-block
density
\[
    q_t(x_t,\alpha,x_{t-1})
    =
    p(x_{t-1},\alpha\mid y_{1:t-1})
    f_\alpha(x_t\mid x_{t-1})
    g_\alpha(y_t\mid x_t).
\]
The next filter is obtained by integrating out \(x_{t-1}\):
\[
    p(x_t,\alpha\mid y_{1:t})
    =
    \frac{\int q_t(x_t,\alpha,x_{t-1})\dd x_{t-1}}
    {\int q_t(x_t,\alpha,x_{t-1})\dd x_t\dd\alpha\dd x_{t-1}}.
\]
In low dimension one can attack this integral by quadrature, particles, or
Gaussian projection.  Zhao--Cui's idea is different: approximate the whole
function \(q_t\) in a separable tensor-train form so that marginal integrals are
cheap contractions instead of high-dimensional numerical integrations.

Implementation object.  A minimal implementation must expose evaluators for
\(\widehat p_{t-1}\), \(f_\alpha\), and \(g_\alpha\), then build an evaluator
for \(q_t\).  The first failure check is that \(q_t\) must be finite and
nonnegative at all fitting points; otherwise neither a square root nor an
importance weight is meaningful.

\section{State-Space Model And BayesFilter Notation}

\paragraph{Source units S1-Eq1--Eq2.}
The first two source equations answer: what conditional independence structure
defines the state-space model?  The variables \(X_t\in\mathcal X\subseteq\R^m\)
are latent states, \(Y_t\in\mathcal Y\subseteq\R^n\) are observations, and
\(\alpha\in\Theta\subseteq\R^d\) is the static parameter.

Zhao--Cui use \(\theta\) for the unknown parameter.  In this note the parameter
is written \(\alpha\), because \(\theta\) is often used later as a generic
argument in derivative formulas.  The paper's model equations (1)--(2) become
\[
    X_t\mid (X_{0:t-1}=x_{0:t-1},\alpha)
    \equiv
    X_t\mid (X_{t-1}=x_{t-1},\alpha)
    \sim f_\alpha(x_t\mid x_{t-1}),
    \tag{BF-1}
\]
and
\[
    Y_t\mid (X_{0:t}=x_{0:t},Y_{1:t-1}=y_{1:t-1},\alpha)
    \equiv
    Y_t\mid (X_t=x_t,\alpha)
    \sim g_\alpha(y_t\mid x_t).
    \tag{BF-2}
\]
The first statement says the latent state process is Markov once \(\alpha\) is
fixed.  The second says the new observation only sees the current state, again
once \(\alpha\) is fixed.

Implementation contract for S1-Eq1--Eq2.  Inputs are \((x_t,x_{t-1},\alpha)\)
for the transition evaluator and \((y_t,x_t,\alpha)\) for the observation
evaluator.  Outputs are nonnegative densities or log densities.  Persisted
state is the model specification and the observed \(y_t\).  Diagnostics are
finite log-density checks and shape checks for \(m,n,d\).

\paragraph{Source unit S1-Eq3.}
The next question is: how do the local transition and observation densities
combine into the full joint density?

The joint density of parameter, states, and observations through time \(t\),
corresponding to Zhao--Cui Eq. (3), is
\[
    p(\alpha,x_{0:t},y_{1:t})
    =
    p(\alpha)p_\alpha(x_0)
    \prod_{j=1}^{t}
    f_\alpha(x_j\mid x_{j-1})g_\alpha(y_j\mid x_j).
    \tag{BF-3}
\]
This is obtained by the chain rule:
\[
\begin{aligned}
 p(\alpha,x_{0:t},y_{1:t})
 &=
 p(\alpha)p_\alpha(x_0)
 \prod_{j=1}^t
 p(x_j\mid x_{0:j-1},y_{1:j-1},\alpha)
 p(y_j\mid x_{0:j},y_{1:j-1},\alpha)\\
 &=
 p(\alpha)p_\alpha(x_0)
 \prod_{j=1}^t
 f_\alpha(x_j\mid x_{j-1})g_\alpha(y_j\mid x_j),
\end{aligned}
\]
where the last equality uses the Markov and observation assumptions.

The derivation is the probabilistic analogue of multiplying conditional
reaction probabilities along a time-ordered chain.  Every new time index adds
one transition factor and one observation factor.  An implementation stores
this as a log-density sum, not as a product, to avoid underflow:
\[
    \log p(\alpha,x_{0:t},y_{1:t})
    =
    \log p(\alpha)+\log p_\alpha(x_0)
    +
    \sum_{j=1}^t
    \{\log f_\alpha(x_j\mid x_{j-1})+\log g_\alpha(y_j\mid x_j)\}.
    \tag{BF-3a}
\]
The diagnostic is simple: any impossible model transition should produce
\(-\infty\) log density intentionally, not a NaN.

\paragraph{Source unit S1-Eq4.}
The fourth source equation answers: how does the joint density become a
posterior density after data are observed?

Conditioning on data gives Zhao--Cui Eq. (4) in BayesFilter notation:
\[
    p(\alpha,x_{0:t}\mid y_{1:t})
    =
    \frac{p(\alpha,x_{0:t},y_{1:t})}
    {p(y_{1:t})},
    \qquad
    p(y_{1:t})
    =
    \int p(\alpha,x_{0:t},y_{1:t})
    \dd\alpha\,\dd x_{0:t}.
    \tag{BF-4}
\]
The denominator is the evidence.  It is a scalar, but it is usually not
available analytically because the integral spans the parameter and all states.

Implementation contract for S1-Eq4.  The numerator is evaluable pointwise.  The
evidence is a scalar integral.  A filtering algorithm that claims a likelihood
must specify how it computes or approximates this scalar.  The failure check is
that the approximated evidence must be positive and finite.

\section{The Four Marginal Learning Problems}

\paragraph{Source units S1-Eq5--Eq8.}
The next four source equations answer: which posterior marginals correspond to
filtering, parameter learning, path estimation, and smoothing?

Zhao--Cui Eqs. (5)--(8) are all marginalizations of the same posterior
\(p(\alpha,x_{0:t}\mid y_{1:t})\).  Filtering keeps only the current state:
\[
    p(x_t\mid y_{1:t})
    =
    \int p(\alpha,x_{0:t}\mid y_{1:t})
    \dd x_{0:t-1}\dd\alpha .
    \tag{BF-5}
\]
Parameter learning keeps only \(\alpha\):
\[
    p(\alpha\mid y_{1:t})
    =
    \int p(\alpha,x_{0:t}\mid y_{1:t})\dd x_{0:t}.
    \tag{BF-6}
\]
Path learning keeps the full state path but removes \(\alpha\):
\[
    p(x_{0:t}\mid y_{1:t})
    =
    \int p(\alpha,x_{0:t}\mid y_{1:t})\dd\alpha .
    \tag{BF-7}
\]
Fixed-lag or retrospective smoothing keeps one old state:
\[
    p(x_k\mid y_{1:t})
    =
    \int p(\alpha,x_{0:t}\mid y_{1:t})
    \dd x_{0:k-1}\dd x_{k+1:t}\dd\alpha ,
    \qquad k<t .
    \tag{BF-8}
\]

The important point is not the names.  It is that the same high-dimensional
posterior density produces all four outputs by integration.  A method that
stores only a cloud of samples can estimate these integrals by sample averages.
A method that stores a Gaussian can compute only Gaussian marginals.  A method
that stores a tensor-train density aims to keep a non-Gaussian function
representation while making many marginal integrals cheap.

Implementation contract for S1-Eq5--Eq8.  The stored object must support
integrating out named coordinate blocks.  The outputs are marginal density
evaluators.  The failure mode is inconsistent normalization: if integrating a
marginal over its retained coordinates does not give one, the contraction or
normalizer is wrong.

\section{The Exact Recursive Bottleneck}

\paragraph{Source units S2.1-Eq9--Eq11.}
These source equations answer: how can the posterior update be performed
sequentially without storing the whole old path forever?

We now derive Zhao--Cui Eqs. (9)--(11).  Start from Bayes' rule:
\[
 p(\alpha,x_{0:t}\mid y_{1:t})
 =
 \frac{p(\alpha,x_{0:t},y_{1:t})}{p(y_{1:t})}.
\]
Using the factorization (BF-3),
\[
 p(\alpha,x_{0:t},y_{1:t})
 =
 p(\alpha,x_{0:t-1},y_{1:t-1})
 f_\alpha(x_t\mid x_{t-1})g_\alpha(y_t\mid x_t).
\]
Also,
\[
 p(\alpha,x_{0:t-1},y_{1:t-1})
 =
 p(\alpha,x_{0:t-1}\mid y_{1:t-1})p(y_{1:t-1}).
\]
Therefore
\[
\begin{aligned}
 p(\alpha,x_{0:t}\mid y_{1:t})
 &=
 \frac{
 p(\alpha,x_{0:t-1}\mid y_{1:t-1})p(y_{1:t-1})
 f_\alpha(x_t\mid x_{t-1})g_\alpha(y_t\mid x_t)}
 {p(y_{1:t})}\\
 &=
 \frac{
 p(\alpha,x_{0:t-1}\mid y_{1:t-1})
 f_\alpha(x_t\mid x_{t-1})g_\alpha(y_t\mid x_t)}
 {p(y_t\mid y_{1:t-1})}.
\end{aligned}
\tag{BF-9}
\]
The last step uses \(p(y_{1:t})=p(y_t\mid y_{1:t-1})p(y_{1:t-1})\).

Now integrate (BF-9) over \(x_{0:t-2}\).  The only factor depending on the old
path \(x_{0:t-2}\) is
\(p(\alpha,x_{0:t-1}\mid y_{1:t-1})\), so
\[
\begin{aligned}
 p(\alpha,x_t,x_{t-1}\mid y_{1:t})
 &=
 \int p(\alpha,x_{0:t}\mid y_{1:t})\dd x_{0:t-2}\\
 &=
 \frac{
 p(\alpha,x_{t-1}\mid y_{1:t-1})
 f_\alpha(x_t\mid x_{t-1})g_\alpha(y_t\mid x_t)}
 {p(y_t\mid y_{1:t-1})}.
\end{aligned}
\tag{BF-10}
\]
Finally, integrate over \(x_{t-1}\) to obtain the next parameter-state filter:
\[
    p(\alpha,x_t\mid y_{1:t})
    =
    \int p(\alpha,x_t,x_{t-1}\mid y_{1:t})\dd x_{t-1}.
    \tag{BF-11}
\]

The numerator in (BF-10) is pointwise evaluable if the previous filter is
evaluable.  The denominator and the marginalization in (BF-11) are the hard
parts.  Zhao--Cui turn this into a function-approximation problem by replacing
the unnormalized numerator with a tensor train.

Implementation contract for S2.1-Eq9--Eq11.  Inputs are the retained evaluator
for \(p(\alpha,x_{t-1}\mid y_{1:t-1})\), the new observation \(y_t\), and model
evaluators \(f_\alpha,g_\alpha\).  The output of the numerator step is a
pointwise evaluator for the three-block target.  The persisted state after
marginalization is an evaluator for \(p(\alpha,x_t\mid y_{1:t})\) and an
evidence increment.  The failure check is that the marginal integral over
\((\alpha,x_t)\) must equal one after normalization.

\section{Tensor Trains From First Principles}

\paragraph{Source units S2.2-TT-display and basis expansion.}
The question answered here is: what separable function class will replace a
full high-dimensional grid?

This section starts from the object a numerical reader can picture: a table of
function values.  If \(D=2\), a function \(h(r_1,r_2)\) sampled on a grid is a
matrix.  Low-rank matrix approximation writes
\[
    h(r_1,r_2)\approx \sum_{\rho=1}^{R} u_\rho(r_1)v_\rho(r_2).
\]
Each term is separable: one factor depends only on \(r_1\), one factor only on
\(r_2\).  A tensor train is the many-variable analogue.  It does not require a
full \(D\)-dimensional table.  It stores a chain of low-rank links, so that
information can pass from coordinate \(1\) to coordinate \(D\) through the rank
indices.

Let \(r=(r_1,\ldots,r_D)\in\Omega_1\times\cdots\times\Omega_D\).  A functional
tensor train approximates a scalar function \(h(r)\) by multiplying
matrix-valued one-dimensional functions:
\[
    h(r)
    \approx
    \widehat h(r)
    =
    H_1(r_1)H_2(r_2)\cdots H_D(r_D),
    \tag{TT-1}
\]
where
\[
    H_k(r_k)\in \R^{R_{k-1}\times R_k},\qquad R_0=R_D=1.
\]
Zhao--Cui's displayed TT formula includes the endpoint rank indices:
\[
    \widehat h(r)
    =
    \sum_{\rho_0=1}^{R_0}\sum_{\rho_1=1}^{R_1}\cdots
    \sum_{\rho_D=1}^{R_D}
    H_1^{(\rho_0,\rho_1)}(r_1)\cdots
    H_D^{(\rho_{D-1},\rho_D)}(r_D),
    \qquad R_0=R_D=1.
    \tag{TT-1a}
\]
The endpoint sums have only one term.  They are written so every core has the
same two rank labels.  A programmer usually evaluates the matrix product
(TT-1); a proof often expands the rank contractions as in (TT-1a).

Because the first object has one row and the last has one column, the product is
a \(1\times 1\) scalar.  Written with indices,
\[
    \widehat h(r)
    =
    \sum_{\rho_1=1}^{R_1}\cdots
    \sum_{\rho_{D-1}=1}^{R_{D-1}}
    H_1^{1,\rho_1}(r_1)
    H_2^{\rho_1,\rho_2}(r_2)\cdots
    H_D^{\rho_{D-1},1}(r_D).
    \tag{TT-2}
\]
The numbers \(R_k\) are TT ranks.  A high rank lets the representation express
strong interactions across a split \((r_1,\ldots,r_k)\) and
\((r_{k+1},\ldots,r_D)\).  A low rank is useful only when the function has
low-dimensional structure across that split.

To see the role of the ranks, split the variables into a left block and a right
block:
\[
    r_{\le k}=(r_1,\ldots,r_k),\qquad
    r_{>k}=(r_{k+1},\ldots,r_D).
\]
Equation (TT-2) can be regrouped as
\[
    \widehat h(r_{\le k},r_{>k})
    =
    \sum_{\rho=1}^{R_k}
    L_\rho(r_{\le k})R_\rho(r_{>k}),
    \tag{TT-2a}
\]
where
\[
    L_\rho(r_{\le k})
    =
    \sum_{\rho_1,\ldots,\rho_{k-1}}
    H_1^{1,\rho_1}(r_1)\cdots H_k^{\rho_{k-1},\rho}(r_k),
\]
and
\[
    R_\rho(r_{>k})
    =
    \sum_{\rho_{k+1},\ldots,\rho_{D-1}}
    H_{k+1}^{\rho,\rho_{k+1}}(r_{k+1})\cdots
    H_D^{\rho_{D-1},1}(r_D).
\]
Thus \(R_k\) is the number of separated terms allowed across that split.  If a
posterior density has only local or weak long-range dependence in the chosen
ordering, moderate ranks can be enough.  If a posterior has strong global
coupling, the ranks grow and the method becomes expensive.

The low-dimensional analogy is the singular value decomposition of a matrix.
For a two-coordinate function, rank \(R\) means \(R\) separable columns times
rows.  For many coordinates, \(R_k\) plays the same role at the cut after
coordinate \(k\).  The failure diagnostic is rank growth: if any fitted
\(R_k\) approaches the imposed maximum, the chosen coordinate ordering, basis,
or preconditioner is not expressing the density compactly.

Each univariate matrix entry is represented in a basis:
\[
    H_k^{a,b}(r_k)
    =
    \sum_{\ell=0}^{p_k-1}
    C_k[a,\ell,b]\psi_{k,\ell}(r_k).
    \tag{TT-3}
\]
This is the paper's basis expansion with renamed coefficient tensor.  Zhao--Cui
write the coefficient as \(A_k[\alpha_{k-1},j,\alpha_k]\); here it is
\(C_k[a,\ell,b]\).  The three axes are left rank, basis index, and right rank.
The basis index describes one-coordinate shape; the rank indices transmit
dependence to neighboring coordinates.

The tensor \(C_k\in\R^{R_{k-1}\times p_k\times R_k}\) is the \(k\)-th TT core.
Thus an implementation stores the basis family, the domain map for each
coordinate, and the core arrays \(C_1,\ldots,C_D\).  To evaluate the function at
one point \(r\), it evaluates all basis vectors
\[
    \psi_k(r_k)
    =
    (\psi_{k,0}(r_k),\ldots,\psi_{k,p_k-1}(r_k))^\top,
\]
forms the matrices \(H_k(r_k)\), and multiplies them.

Implementation contract for S2.2.  Inputs are a coordinate \(r_k\), basis
metadata, and core tensor \(C_k\).  The operation is basis evaluation followed
by a contraction over the basis index:
\[
    H_k(r_k)=\sum_{\ell=0}^{p_k-1}\psi_{k,\ell}(r_k)\,C_k[:,\ell,:].
    \tag{TT-3a}
\]
The persisted state is the list of cores, basis objects, and coordinate
domains.  The failure checks are basis overflow, invalid coordinate outside the
domain, and incompatible adjacent ranks.

\subsection{A Concrete Basis Choice}

\paragraph{Source unit S2.2-basis-family.}
Zhao--Cui allow several basis families.  This note fixes one basis to make the
paper's abstract core formula executable.

Zhao--Cui's software supports several basis families.  For a self-contained
minimal implementation, choose one.  On a finite interval
\([a_k,b_k]\), map a physical coordinate \(r_k\) to
\[
    z_k=\frac{2r_k-(a_k+b_k)}{b_k-a_k}\in[-1,1],
    \qquad
    r_k=\frac{a_k+b_k}{2}+\frac{b_k-a_k}{2}z_k.
    \tag{TT-6}
\]
Use normalized Legendre polynomials
\[
    \psi_{k,\ell}(r_k)
    =
    \sqrt{\frac{2\ell+1}{b_k-a_k}}\,
    P_\ell(z_k),
    \qquad \ell=0,\ldots,p_k-1.
    \tag{TT-7}
\]
The Legendre polynomials are evaluated by the recurrence
\[
    P_0(z)=1,\qquad P_1(z)=z,
\]
\[
    (\ell+1)P_{\ell+1}(z)
    =
    (2\ell+1)zP_\ell(z)-\ell P_{\ell-1}(z),
    \qquad \ell\ge1.
    \tag{TT-8}
\]
This basis is useful pedagogically because its mass matrix is the identity:
\[
    \int_{a_k}^{b_k}\psi_{k,\ell}(r_k)\psi_{k,\ell'}(r_k)\dd r_k
    =
    \delta_{\ell\ell'}.
    \tag{TT-9}
\]
If another basis is used, every formula below remains the same after replacing
the identity mass matrix by the corresponding one-dimensional mass matrix.

Zhao--Cui state that functional alternating schemes with cross interpolation
typically use
\[
    O(DpR^2)\quad\text{target evaluations}
    \tag{TT-10}
\]
and
\[
    O(DpR^3)\quad\text{floating point operations},
    \tag{TT-11}
\]
where \(p=\max_k p_k\) and \(R=\max_k R_k\).  This is a cost statement after
the ranks are moderate; it is not a promise that \(R\) stays small.  The
implementation should therefore record fitted ranks, residuals, and
ill-conditioning diagnostics.

\section{Why TT Marginalization Is Cheap Once The Representation Exists}

\paragraph{Source unit S2.2-core-integration.}
The question answered here is: once a function is in TT form, how does the
algorithm integrate out one coordinate without returning to a full grid?

Suppose \(h\) is represented by (TT-1).  If we need to integrate out coordinate
\(r_k\), the product structure gives
\[
\begin{aligned}
 \int \widehat h(r_1,\ldots,r_D)\dd r_k
 &=
 H_1(r_1)\cdots H_{k-1}(r_{k-1})
 \left(\int H_k(r_k)\dd r_k\right)
 H_{k+1}(r_{k+1})\cdots H_D(r_D).
\end{aligned}
\tag{TT-4}
\]
Only the \(k\)-th core changes.  Its integrated matrix is
\[
    \overline H_k
    =
    \int H_k(r_k)\dd r_k,\qquad
    \overline H_k[a,b]
    =
    \sum_{\ell=0}^{p_k-1} C_k[a,\ell,b]
    \int \psi_{k,\ell}(r_k)\dd r_k.
    \tag{TT-5}
\]
If the one-dimensional basis integrals are precomputed, marginalizing one
coordinate is a small matrix contraction.  Repeatedly applying (TT-4) integrates
whole blocks of variables.

This explains why Zhao--Cui target a TT representation of
\((x_t,\alpha,x_{t-1})\).  Once that three-block function is separable enough,
the new filter \(p(x_t,\alpha\mid y_{1:t})\) is obtained by integrating the
old-state block \(x_{t-1}\), one coordinate at a time.

Implementation contract for S2.2-core-integration.  Inputs are the TT cores and
precomputed one-dimensional basis integrals.  The operation replaces one core
by \(\overline H_k\) and contracts adjacent matrices.  The persisted object is a
lower-dimensional TT or evaluator.  The failure check is dimension/order
consistency: integrating the wrong block silently returns the wrong filter.

\section{Zhao--Cui Algorithm 1, Fully Annotated}

\paragraph{Source unit S2.3-Algorithm1(a).}
Algorithm 1(a) answers: what target can be evaluated at time \(t\) using only
the retained approximation from time \(t-1\)?

The exact recursion (BF-10) contains the previous exact filter.  A numerical
algorithm has the previous approximation.  Let
\(\widehat\pi_{t-1}(x_{t-1},\alpha)\) denote an unnormalized approximation of
\(p(x_{t-1},\alpha\mid y_{1:t-1})\).  Zhao--Cui Eq. (12) becomes
\[
    q_t(x_t,\alpha,x_{t-1})
    =
    \widehat\pi_{t-1}(x_{t-1},\alpha)
    f_\alpha(x_t\mid x_{t-1})
    g_\alpha(y_t\mid x_t).
    \tag{A1-1}
\]
This is not separable in general.  Even if
\(\widehat\pi_{t-1}\) is a TT over \((\alpha,x_{t-1})\), the transition and
likelihood couple the new and old state coordinates.

Zhao--Cui use approximate proportionality before normalization.  In explicit
terms, \(h_1\propto_{\rm approx} h_2\) means
\[
    \frac{h_1(r)}{\int h_1(s)\dd s}
    \approx
    \frac{h_2(r)}{\int h_2(s)\dd s}.
    \tag{A1-0}
\]
This is why the algorithm can store an unnormalized TT.  The normalization is
not forgotten; it is postponed until the complete TT contraction
\(\widehat Z_t\) is computed.

Algorithm 1 has three mathematical steps.

\paragraph{Step 1: build a pointwise target.}
Given \(r_t=(x_t,\alpha,x_{t-1})\), evaluate \(q_t(r_t)\) by multiplying the
previous approximation, transition density, and likelihood:
\[
    r_t\mapsto
    \widehat\pi_{t-1}(x_{t-1},\alpha)
    f_\alpha(x_t\mid x_{t-1})
    g_\alpha(y_t\mid x_t).
\]
Storage needed: an evaluator for the previous TT approximation and evaluators
for \(f_\alpha\) and \(g_\alpha\).

Mini implementation contract.  Inputs:
\(\widehat\pi_{t-1}\), \(f_\alpha\), \(g_\alpha\), \(y_t\), and a point
\((x_t,\alpha,x_{t-1})\).  Output: \(q_t\) or \(\log q_t\).  Persisted state:
none beyond the evaluator.  Operation: multiply three factors or add three log
factors.  Diagnostic: reject NaN, negative density values, or inconsistent
coordinate block sizes.

\paragraph{Step 2: re-approximate by a TT.}
\paragraph{Source unit S2.3-Algorithm1(b).}
This source unit answers: how does the nonseparable target become a separable
object that can be marginalized?

Choose a variable order, in the paper usually
\[
    (x_t,\alpha,x_{t-1}).
\]
Fit a TT
\[
    \widehat\pi_t(x_t,\alpha,x_{t-1})
    =
    F_{t,1}(x_{t,1})\cdots F_{t,m_x}(x_{t,m_x})
    A_{t,1}(\alpha_1)\cdots A_{t,m_\alpha}(\alpha_{m_\alpha})
    H_{t,1}(x_{t-1,1})\cdots H_{t,m_x}(x_{t-1,m_x})
    \tag{A1-2}
\]
so that \(\widehat\pi_t\approx q_t\).  Zhao--Cui use functional TT tools with
alternating approximation and rank adaptation.  For implementation one must
choose a domain map, basis, fitting points, rank rule, and stopping rule.  The
paper's research code supplies one such family; the fixed-branch section below
supplies a differentiable BayesFilter specialization.

The most direct implementation view is least squares.  Choose fitting points
\[
    r^{(j)}=(x_t^{(j)},\alpha^{(j)},x_{t-1}^{(j)}),
    \qquad j=1,\ldots,N_{\rm fit},
\]
and define target values
\[
    y_j=q_t(r^{(j)}).
\]
For fixed TT ranks and fixed all cores except core \(k\), the fitted function is
linear in the entries of core \(C_k\).  There is a design row \(A_{j,k}\) such
that
\[
    \widehat\pi_t(r^{(j)})=A_{j,k}g_k,
    \qquad g_k=\vecop(C_k).
    \tag{A1-2a}
\]
Let \(L_{j,k}\) be the left environment obtained by multiplying all cores before
\(k\) at point \(r^{(j)}\), and let \(R_{j,k}\) be the right environment from
all cores after \(k\).  Let
\(\psi_k(r_k^{(j)})\) be the local basis vector.  Then
\[
    A_{j,k}
    =
    R_{j,k}^{\top}\otimes \psi_k(r_k^{(j)})^\top\otimes L_{j,k}.
    \tag{A1-2b}
\]
This formula is worth lingering over.  The left environment tells the current
core what the earlier coordinates have already contributed.  The right
environment tells it what later coordinates will contribute.  The basis vector
gives the local coordinate dependence.  The Kronecker product stitches these
three pieces into a row that multiplies the vectorized core.

With weights \(w_j>0\), the core update solves
\[
    \min_{g_k}
    \sum_{j=1}^{N_{\rm fit}}w_j\bigl(A_{j,k}g_k-y_j\bigr)^2
    +\rho\|g_k\|_2^2,
    \tag{A1-2c}
\]
or equivalently
\[
    (A_k^\top W A_k+\rho I)g_k=A_k^\top Wy.
    \tag{A1-2d}
\]
An alternating scheme sweeps through the cores, updating one core at a time.
Zhao--Cui's adaptive implementation is more sophisticated, but this least-
squares view is the easiest way to see how a TT approximation becomes an
actual numerical object.

Mini implementation contract.  Inputs: target evaluator \(q_t\), domains,
basis objects, fitting points, weights, ranks, and sweep settings.  Output:
cores \(F,A,H\) for the fitted TT.  Persisted state: cores, rank vector, basis
metadata, fit residual, and any scaling shift.  Operation: alternating core
least-squares or cross approximation.  Diagnostic: weighted residual, condition
number of each normal matrix, rank growth, and negative values if this basic
non-squared algorithm is interpreted as a density.

\paragraph{Step 3: integrate the old state.}
\paragraph{Source unit S2.3-Algorithm1(c).}
This source unit answers: after the three-block target is fitted, what is
retained for the next sequential update?

The next approximate filter is
\[
    \widehat\pi_t(x_t,\alpha)
    =
    \int \widehat\pi_t(x_t,\alpha,x_{t-1})\dd x_{t-1}.
    \tag{A1-3}
\]
Because \(x_{t-1}\) occupies the last block, this integral is a sequence of core
integrals:
\[
    \widehat\pi_t(x_t,\alpha)
    =
    F_{t,1}(x_{t,1})\cdots F_{t,m_x}(x_{t,m_x})
    A_{t,1}(\alpha_1)\cdots A_{t,m_\alpha}(\alpha_{m_\alpha})
    \overline H_{t,1}\cdots \overline H_{t,m_x}.
    \tag{A1-4}
\]
The approximate normalizer is
\[
    \widehat Z_t
    =
    \int \widehat\pi_t(x_t,\alpha,x_{t-1})
    \dd x_t\,\dd\alpha\,\dd x_{t-1},
    \tag{A1-5}
\]
again a complete TT contraction.

The same fitted TT gives the parameter marginal and filtering marginal:
\[
    \widehat p_t(\alpha)
    =
    \frac{\int\widehat\pi_t(x_t,\alpha,x_{t-1})
    \dd x_t\,\dd x_{t-1}}
    {\widehat Z_t},
    \qquad
    \widehat p_t(x_t)
    =
    \frac{\int\widehat\pi_t(x_t,\alpha,x_{t-1})
    \dd\alpha\,\dd x_{t-1}}
    {\widehat Z_t}.
    \tag{A1-6}
\]
The retained object for the next time step is not merely a scalar
normalizer.  It is the evaluable two-block density
\[
    \widehat\pi_t(x_t,\alpha)
    =
    \int\widehat\pi_t(x_t,\alpha,x_{t-1})\dd x_{t-1},
    \tag{A1-7}
\]
because time \(t+1\) will plug this function into the next target
\(q_{t+1}\).

Mini implementation contract.  Inputs: fitted three-block TT and block ordering
\((x_t,\alpha,x_{t-1})\).  Outputs: retained evaluator
\(\widehat\pi_t(x_t,\alpha)\), normalizer \(\widehat Z_t\), optional marginals
\(\widehat p_t(x_t)\) and \(\widehat p_t(\alpha)\).  Persisted state: retained
TT/evaluator and contraction metadata.  Operation: endpoint TT core
integration.  Diagnostic: \(\widehat Z_t\) must be finite and positive, and
normalizing retained marginals must integrate to one.

The main failure mode of Algorithm 1 is sign.  A fitted TT is a real-valued
function approximation.  Truncation or interpolation can make
\(\widehat\pi_t(r_t)<0\), which is impossible for a density.  Negative values
also break transport maps and importance-sampling ratios.  This motivates the
squared-TT construction.

\section{Why Nonnegativity Fails And Why Square-Root TT Repairs It}

\paragraph{Source unit S3.1-nonnegativity-barrier.}
This source unit answers: why is Algorithm 1 not enough for transport maps and
importance correction?

Approximating a nonnegative function by a truncated expansion does not preserve
nonnegativity.  The one-dimensional analogy is simple:
\[
    h(x)\ge 0
    \quad\text{but}\quad
    \widehat h(x)=\sum_{\ell=0}^{p-1}c_\ell\psi_\ell(x)
    \text{ may be negative.}
\]
The tensor-train case has the same issue.  A low-rank approximation can reduce
error in a least-squares sense while still creating small negative regions.

Zhao--Cui's repair is to approximate the square root of an unnormalized density.
If
\[
    \sqrt{\pi(r)}\approx \phi(r),
\]
then
\[
    \phi(r)^2\ge 0
\]
for all \(r\).  A small defensive density is then added so that the proposal
density is positive wherever the target might be positive.

\section{Squared-TT Density, Defensive Reference, And Normalizer}

\paragraph{Source units S3.1-Eq13 and Lemma 1.}
The question answered here is: how can a TT approximation be forced to define a
nonnegative density whose support covers the target?

Zhao--Cui Eq. (13) becomes the following.  Let \(\pi(r)\ge0\) be an
unnormalized target on a domain \(\Omega\subseteq\R^D\), with normalizer
\[
    Z=\int_\Omega \pi(r)\dd r.
\]
Let \(\lambda(r)\) be a normalized reference density on \(\Omega\),
\[
    \lambda(r)\ge0,\qquad \int_\Omega \lambda(r)\dd r=1,
\]
and choose \(\tau>0\).  Fit a TT \(\phi(r)\) to \(\sqrt{\pi(r)}\).  Define
\[
    \widehat\pi(r)
    =
    \phi(r)^2+\tau\lambda(r),
    \qquad
    \widehat Z
    =
    \int_\Omega \widehat\pi(r)\dd r,
    \qquad
    \widehat p(r)
    =
    \frac{\widehat\pi(r)}{\widehat Z}.
    \tag{S1}
\]
Since \(\lambda\) is normalized,
\[
    \widehat Z
    =
    \int_\Omega \phi(r)^2\dd r+\tau.
    \tag{S2}
\]
The defensive term has two jobs.  First, it keeps the approximation positive
even if \(\phi\) is zero somewhere.  Second, if \(\pi/\lambda\) is bounded, then
\[
    \frac{p(r)}{\widehat p(r)}
    =
    \frac{\pi(r)/Z}{\widehat\pi(r)/\widehat Z}
    =
    \frac{\widehat Z}{Z}
    \frac{\pi(r)}{\phi(r)^2+\tau\lambda(r)}
    \le
    \frac{\widehat Z}{Z\tau}
    \frac{\pi(r)}{\lambda(r)}.
    \tag{S3}
\]
This is the importance-sampling reason for the defensive density: the true
target is absolutely continuous with respect to the approximation when the
reference covers the target support.

Zhao--Cui's Lemma 1 ties the square-root fit to density error.  In the notation
above, suppose
\[
    \|\phi-\sqrt{\pi}\|_{L^2}\le \varepsilon,
    \qquad
    \sqrt{\tau}\le \|\phi-\sqrt{\pi}\|_{L^2}.
    \tag{S4}
\]
Then the square-root error of the defensive approximation satisfies
\[
    \|\sqrt{\widehat\pi}-\sqrt{\pi}\|_{L^2}\le 2\varepsilon,
    \tag{S5}
\]
and the normalized Hellinger distance obeys
\[
    D_H(\widehat p,p)
    \le
    \frac{\sqrt{2}}{\sqrt Z}
    \|\sqrt{\widehat\pi}-\sqrt{\pi}\|_{L^2}
    \le
    \frac{2\sqrt{2}\varepsilon}{\sqrt Z}.
    \tag{S6}
\]
For this note, the important lesson is local: fitting the square root in
\(L^2\) controls the normalized density in Hellinger distance after
normalization.  The lemma does not say ranks will be small, nor does it make an
adaptive implementation globally differentiable.

\begin{center}
\fbox{\begin{minipage}{0.92\linewidth}
\textbf{Local assumption box for Lemma 1.}
P18 uses only the following local content: if the square-root fit is small in
\(L^2\), then the normalized defensive density is close in Hellinger distance
under the stated support and finite-normalizer conditions.  P18 does not
re-prove the external Cui--Dolgov theorem, does not claim a rank bound, and does
not claim adaptive differentiability.
\end{minipage}}
\end{center}

Implementation contract for S3.1-Eq13.  Inputs are a square-root target
evaluator, a reference density \(\lambda\), a defensive mass \(\tau\), and a
fitted TT \(\phi\).  Outputs are \(\widehat\pi\), \(\widehat Z\), and
\(\widehat p\).  Persisted state is \(\phi\)'s cores, \(\lambda\), \(\tau\),
and \(\widehat Z\).  Diagnostics are positivity, finite \(\widehat Z\), support
coverage by \(\lambda\), and the defensive ratio \(\tau/\widehat Z\).

\section{Squared-TT Marginalization And Mass Matrices}

\paragraph{Source unit S3.1-Proposition2.}
This proposition answers: how do we marginalize a squared TT without expanding
the square into a full tensor?

The square changes the marginalization algebra.  Let
\[
    \phi(r)=H_1(r_1)\cdots H_D(r_D),
    \qquad H_k(r_k)\in\R^{R_{k-1}\times R_k}.
\]
For a split after coordinate \(k\), define
\[
    H_{\le k}(r_{\le k})=H_1(r_1)\cdots H_k(r_k)\in\R^{1\times R_k},
\]
and
\[
    H_{>k}(r_{>k})=H_{k+1}(r_{k+1})\cdots H_D(r_D)\in\R^{R_k\times 1}.
\]
Then
\[
    \phi(r)=H_{\le k}(r_{\le k})H_{>k}(r_{>k}).
\]
The marginal square term is
\[
\begin{aligned}
 \int \phi(r)^2\dd r_{>k}
 &=
 \int
 \left[
 H_{\le k}(r_{\le k})H_{>k}(r_{>k})
 \right]^2\dd r_{>k}\\
 &=
 \sum_{a=1}^{R_k}\sum_{b=1}^{R_k}
 H_{\le k,a}(r_{\le k})
 \left[
 \int H_{>k,a}(r_{>k})H_{>k,b}(r_{>k})\dd r_{>k}
 \right]
 H_{\le k,b}(r_{\le k}).
\end{aligned}
\tag{M1}
\]
Define the right mass matrix
\[
    M_{>k}[a,b]
    =
    \int H_{>k,a}(r_{>k})H_{>k,b}(r_{>k})\dd r_{>k}.
    \tag{M2}
\]
This matrix is positive semidefinite.  If \(M_{>k}=L_{>k}L_{>k}^\top\) is a
Cholesky or square-root factorization, then
\[
    \int \phi(r)^2\dd r_{>k}
    =
    \sum_{\gamma=1}^{R_k}
    \left[
    \sum_{a=1}^{R_k}H_{\le k,a}(r_{\le k})L_{>k}[a,\gamma]
    \right]^2.
    \tag{M3}
\]
Including the defensive reference, if
\(\lambda(r)=\prod_{j=1}^D\lambda_j(r_j)\), gives
\[
    \widehat p(r_{\le k})
    =
    \frac{
    \sum_{\gamma=1}^{R_k}
    \left[
    \sum_{a=1}^{R_k}H_{\le k,a}(r_{\le k})L_{>k}[a,\gamma]
    \right]^2
    +
    \tau\prod_{j=1}^k\lambda_j(r_j)}
    {\widehat Z}.
    \tag{M4}
\]
This is Zhao--Cui Eq. (14) in BayesFilter notation.

The practical recursion starts from the right end.  Set \(M_{>D}=1\).  Moving
from \(j=D\) down to \(1\), compute
\[
    M_{>j-1}
    =
    \int H_j(r_j)M_{>j}H_j(r_j)^\top\dd r_j.
    \tag{M5}
\]
In basis coefficients this integral is finite-dimensional.  If
\[
    H_j^{a,b}(r_j)=\sum_{\ell}C_j[a,\ell,b]\psi_{j,\ell}(r_j),
\]
and
\[
    B_j[\ell,\ell']
    =
    \int \psi_{j,\ell}(r_j)\psi_{j,\ell'}(r_j)\dd r_j,
\]
then
\[
    M_{>j-1}[a,a']
    =
    \sum_{b,b'=1}^{R_j}
    \sum_{\ell,\ell'}
    C_j[a,\ell,b]\,M_{>j}[b,b']\,C_j[a',\ell',b']\,B_j[\ell,\ell'].
    \tag{M6}
\]
This formula is the bridge from the paper to code: marginalization stores mass
matrices or their square roots and never enumerates a full tensor grid.

Implementation contract for S3.1-Proposition2.  Inputs are TT cores, basis mass
matrices \(B_j\), defensive reference factors \(\lambda_j\), and \(\tau\).
Outputs are marginal density evaluators and mass contractions \(M_{>k}\) or
their Cholesky factors.  Persisted state is the list of mass matrices needed by
later conditional CDFs.  Diagnostics are symmetry error
\(\|M-M^\top\|\), Cholesky failure, negative eigenvalues beyond tolerance, and
normalizer mismatch between independent contraction paths.

\begin{center}
\fbox{\begin{minipage}{0.92\linewidth}
\textbf{Local assumption box for Proposition 2.}
P18 uses the finite-dimensional contraction identity for a squared TT with
integrable basis products and positive semidefinite mass matrices.  The note
derives the contraction algebra explicitly in (M1)--(M9).  It does not claim
that every numerical Cholesky factorization is stable; that is a diagnostic
obligation.
\end{minipage}}
\end{center}

\subsection{A Small Three-Coordinate Marginalization}

Take \(D=3\), and suppose
\[
    \phi(r_1,r_2,r_3)=H_1(r_1)H_2(r_2)H_3(r_3).
\]
If \(H_1\in\R^{1\times R_1}\),
\(H_2\in\R^{R_1\times R_2}\), and
\(H_3\in\R^{R_2\times 1}\), then marginalizing out \(r_3\) gives
\[
\begin{aligned}
 \int \phi(r_1,r_2,r_3)^2\dd r_3
 &=
 H_1(r_1)H_2(r_2)
 \left[\int H_3(r_3)H_3(r_3)^\top\dd r_3\right]
 H_2(r_2)^\top H_1(r_1)^\top.
\end{aligned}
\tag{M7}
\]
The bracketed quantity is an \(R_2\times R_2\) matrix.  It is not a function of
\(r_1\) or \(r_2\); it is a stored contraction.  If we also marginalize out
\(r_2\), then
\[
    M_{>1}
    =
    \int
    H_2(r_2)
    \left[\int H_3(r_3)H_3(r_3)^\top\dd r_3\right]
    H_2(r_2)^\top
    \dd r_2,
    \tag{M8}
\]
and the one-coordinate marginal is
\[
    \int\phi(r_1,r_2,r_3)^2\dd r_2\dd r_3
    =
    H_1(r_1)M_{>1}H_1(r_1)^\top.
    \tag{M9}
\]
This is exactly what the recursive mass formula (M5) does.  The notation in the
paper is compact because it assumes the reader already sees this matrix chain.
The implementation sees it as repeated multiplication and integration of small
matrices.

\section{Conditional Densities And KR Maps From Marginal Ratios}

\paragraph{Source units S3.1-KR-ratios and Remark 3.}
The question answered here is: how do marginal densities become a transport map
that can sample from the approximate density?

For a normalized density \(\widehat p(r_1,\ldots,r_D)\), the chain rule gives
\[
    \widehat p(r)
    =
    \widehat p(r_1)
    \prod_{k=2}^D
    \widehat p(r_k\mid r_{<k}),
    \qquad
    \widehat p(r_k\mid r_{<k})
    =
    \frac{\widehat p(r_{\le k})}{\widehat p(r_{<k})}.
    \tag{K1}
\]
The squared-TT marginal formula (M4) supplies both numerator and denominator.
Thus each conditional density is computable from TT contractions.

Define conditional distribution functions
\[
    F_1(r_1)=\int_{-\infty}^{r_1}\widehat p(s_1)\dd s_1,
    \tag{K2}
\]
and, for \(k\ge2\),
\[
    F_k(r_k\mid r_{<k})
    =
    \int_{-\infty}^{r_k}
    \widehat p(s_k\mid r_{<k})\dd s_k.
    \tag{K3}
\]
The lower triangular Knothe--Rosenblatt map is
\[
    F(r)
    =
    \bigl(
    F_1(r_1),
    F_2(r_2\mid r_1),
    \ldots,
    F_D(r_D\mid r_{<D})
    \bigr).
    \tag{K4}
\]
If \(R\sim\widehat p\), then \(F(R)\) is uniform on \([0,1]^D\).  Conversely, if
\(\Xi\sim \mathrm{Unif}([0,1]^D)\), then \(F^{-1}(\Xi)\sim\widehat p\).

The proof is the triangular Jacobian calculation.  The Jacobian of \(F\) is
lower triangular, and
\[
    \frac{\partial F_k(r_k\mid r_{<k})}{\partial r_k}
    =
    \widehat p(r_k\mid r_{<k}).
\]
Therefore
\[
    |\det\nabla F(r)|
    =
    \prod_{k=1}^D \widehat p(r_k\mid r_{<k})
    =
    \widehat p(r).
    \tag{K5}
\]
The pullback of the uniform density by \(F\) is exactly
\(|\det\nabla F(r)|=\widehat p(r)\).

\paragraph{Source unit S3.1-Remark3.}
Remark 3 answers: what changes if the conditional map is built from the right
end rather than from the left end?

The paper also states the reverse-order construction.  If instead we compute
right-to-left marginals
\[
    \widehat p(r_{\ge k})
    =
    \int \widehat p(r)\dd r_{<k},
    \tag{K6}
\]
then the conditional densities are
\[
    \widehat p(r_k\mid r_{>k})
    =
    \frac{\widehat p(r_{\ge k})}{\widehat p(r_{>k})},
    \tag{K7}
\]
and the upper triangular KR map is
\[
    F^u(r)
    =
    \bigl(
    F^u_1(r_1\mid r_{>1}),\ldots,
    F^u_{D-1}(r_{D-1}\mid r_D),F^u_D(r_D)
    \bigr).
    \tag{K8}
\]
This is not a different density approximation.  It is the same density
integrated in the opposite direction.  Zhao--Cui use the lower map for
backward smoothing and the upper map for forward particle proposals.

Implementation contract for Remark 3.  Inputs are the same squared-TT density
and the choice of reverse integration direction.  Outputs are right-to-left
mass contractions and upper-triangular conditional CDFs.  Persisted state must
record the direction, because using lower-map contractions inside an upper-map
proposal silently changes the proposal density.  The failure check is endpoint
normalization of every reverse conditional CDF.

The computational costs in the paper separate preprocessing from per-sample
map evaluation.  Building all conditional densities from the squared-TT mass
recursions costs
\[
    O(DpR^3).
    \tag{K9}
\]
For \(N\) samples, evaluating conditional densities along a triangular map
costs approximately
\[
    O(NDpR^2),
    \tag{K10}
\]
and constructing/evaluating one-dimensional distribution functions adds basis
and root-finding work.  If \(c_{\rm root}\) is the fixed number of root-finding
iterations, inverse-map evaluation scales as
\[
    O(DpR^3+NDpR^2+NDp(\log p+R)+NDp\,c_{\rm root}).
    \tag{K11}
\]
The formula is less important than the decomposition: build the conditional
objects once, then pay a per-sample triangular evaluation cost.

Implementation contract for S3.1-KR-ratios.  Inputs are marginal density
evaluators from Proposition 2.  Outputs are conditional CDF evaluators and
inverse-CDF samplers.  Persisted state includes integration direction, mass
matrices, one-dimensional CDF coefficients, and root-finding tolerances.
Diagnostics are monotonicity of each CDF, endpoint values near \(0\) and \(1\),
and inverse residuals.

\begin{center}
\fbox{\begin{minipage}{0.92\linewidth}
\textbf{Local assumption box for KR exactness.}
P18 uses the triangular CDF identity under positive conditional densities,
well-defined marginal ratios, and invertible one-dimensional conditional CDFs.
If a conditional CDF is flat, discontinuous, numerically non-monotone, or
unsupported by the reference domain, the KR sampler is not valid until that
failure is fixed.
\end{minipage}}
\end{center}

\section{Zhao--Cui Algorithm 2, Fully Annotated}

\paragraph{Source units S3.2-Algorithm2(a)--(c), Eq15--Eq16.}
Algorithm 2 answers: how does the sequential recursion change when the fitted
object is the square root and the stored density is squared?

Algorithm 2 is Algorithm 1 with a square-root TT and a defensive density.
At time \(t-1\), assume the previous filter approximation is evaluable as
\[
    \widehat p_{t-1}(x_{t-1},\alpha).
\]
The pointwise unnormalized target, Zhao--Cui Eq. (15), is
\[
    q_t(x_t,\alpha,x_{t-1})
    =
    \widehat p_{t-1}(x_{t-1},\alpha)
    f_\alpha(x_t\mid x_{t-1})
    g_\alpha(y_t\mid x_t).
    \tag{A2-1}
\]
Fit a TT to the square root:
\[
    \sqrt{q_t(x_t,\alpha,x_{t-1})}
    \approx
    \phi_t(x_t,\alpha,x_{t-1}).
    \tag{A2-2}
\]
Then define Zhao--Cui Eq. (16):
\[
    \widehat q_t(x_t,\alpha,x_{t-1})
    =
    \phi_t(x_t,\alpha,x_{t-1})^2
    +
    \tau_t\lambda_t(x_t)\lambda_\alpha(\alpha)\lambda_{t-1}(x_{t-1}).
    \tag{A2-3}
\]
The approximate normalizer is
\[
    \widehat Z_t
    =
    \int \widehat q_t(x_t,\alpha,x_{t-1})
    \dd x_t\,\dd\alpha\,\dd x_{t-1}.
    \tag{A2-4}
\]
The next normalized approximate filter is
\[
    \widehat p_t(x_t,\alpha)
    =
    \frac{\int \widehat q_t(x_t,\alpha,x_{t-1})\dd x_{t-1}}
    {\widehat Z_t}.
    \tag{A2-5}
\]

Everything is now nonnegative.  The approximation risk has moved: the method
must fit the square root accurately enough, choose a defensive term that is not
too small for stability or too large for accuracy, and keep TT ranks manageable.

Mini implementation contract.  Inputs: previous retained density evaluator,
model evaluators, reference factors \(\lambda_t,\lambda_\alpha,\lambda_{t-1}\),
\(\tau_t\), basis/rank/fitting controls, and observation \(y_t\).  Outputs:
square-root TT \(\phi_t\), nonnegative density evaluator \(\widehat q_t\),
normalizer \(\widehat Z_t\), retained filter \(\widehat p_t(x_t,\alpha)\), and
mass contractions for conditional maps.  Persisted state: all of those objects.
Diagnostics: square-root fit residual, defensive mass ratio, \(\widehat Z_t>0\),
rank growth, and Cholesky/mass-matrix health.

\section{Forward Conditional Map And Particle-Filter Correction}

\paragraph{Source units S3.3-Eq20--Eq23 and Algorithm3.}
These source units answer: how does the squared-TT approximation become a
particle proposal, and how is approximation bias corrected?

The same \(\widehat q_t(x_t,\alpha,x_{t-1})\) defines a conditional density for
the new state given the old state and parameter:
\[
    \widehat p_t(x_t\mid \alpha,x_{t-1},y_{1:t})
    =
    \frac{\widehat q_t(x_t,\alpha,x_{t-1})}
    {\int \widehat q_t(s,\alpha,x_{t-1})\dd s}.
    \tag{F1}
\]
Constructing the corresponding upper triangular conditional KR map gives a
proposal sampler:
\[
    \widehat X_t^{(i)}
    =
    (F_{t,t}^{u})^{-1}
    \left(\Xi_t^{(i)}\mid \widehat\alpha^{(i)},\widehat x_{t-1}^{(i)}\right),
    \qquad
    \Xi_t^{(i)}\sim\mathrm{Unif}([0,1]^{m_x}).
    \tag{F2}
\]
This is Zhao--Cui Eq. (21) in notation adapted to BayesFilter.

The proposal density for the full path after sampling is
\[
    \widehat p_t(x_t\mid \alpha,x_{t-1},y_{1:t})
    p(\alpha,x_{0:t-1}\mid y_{1:t-1}),
    \tag{F3}
\]
which corresponds to Zhao--Cui Eq. (22).  The exact posterior recursion is
\[
    p(\alpha,x_{0:t}\mid y_{1:t})
    \propto
    p(\alpha,x_{0:t-1}\mid y_{1:t-1})
    f_\alpha(x_t\mid x_{t-1})
    g_\alpha(y_t\mid x_t).
\]
Dividing exact target by proposal gives the unnormalized correction factor
\[
    \omega_f(\alpha,x_{t-1:t})
    =
    \frac{
    f_\alpha(x_t\mid x_{t-1})g_\alpha(y_t\mid x_t)}
    {\widehat p_t(x_t\mid \alpha,x_{t-1},y_{1:t})}.
    \tag{F4}
\]
This is Zhao--Cui Eq. (23).  The correction is conceptually important: the TT
proposal can be biased, but if its conditional density is evaluable and covers
the target, importance weights can correct expectations.

With incoming normalized weights \(W_{t-1}^{(i)}\), the particle-filter weight
update is
\[
    \widetilde W_t^{(i)}
    =
    W_{t-1}^{(i)}
    \omega_f(\widehat\alpha^{(i)},\widehat x_{t-1:t}^{(i)}),
    \tag{F5}
\]
followed by
\[
    W_t^{(i)}
    =
    \frac{\widetilde W_t^{(i)}}
    {\sum_{j=1}^{N}\widetilde W_t^{(j)}}.
    \tag{F6}
\]
Algorithm 3 is therefore an ordinary importance sampler whose proposal is
defined by the upper conditional KR map of the TT approximation.  The
implementation stores the conditional-map object, the sampled \(x_t^{(i)}\),
the proposal density evaluator, and the normalized weights.  A standard failure
diagnostic is effective sample size,
\[
    \operatorname{ESS}_t
    =
    \frac{1}{\sum_{i=1}^{N}(W_t^{(i)})^2}.
    \tag{F7}
\]
Low ESS means that the TT proposal is not close enough to the exact recursive
posterior for stable correction.

Mini implementation contract.  Inputs: weighted particles at time \(t-1\), the
upper conditional map, model density evaluators, and proposal density evaluator.
Outputs: new particles and normalized weights.  Persisted state: particles,
weights, ESS, and optional resampling history.  Operation: inverse triangular
sampling plus importance-weight update.  Diagnostics: ESS, infinite/zero
proposal density, root-finding failure, and weight underflow.

\section{Backward Conditional Map, Path Estimation, And Smoothing}

\paragraph{Source units S3.4-Eq24--Eq26, Figure 1, and Algorithm4.}
These source units answer: how can the same sequence of local TT approximations
produce full paths and smoothed old states after all data are observed?

The lower conditional map goes the other direction:
\[
    \widehat p_t(x_{t-1}\mid x_t,\alpha,y_{1:t})
    =
    \frac{\widehat q_t(x_t,\alpha,x_{t-1})}
    {\int \widehat q_t(x_t,\alpha,s)\dd s}.
    \tag{B1}
\]
The corresponding inverse conditional KR map samples
\[
    \widehat X_{t-1}^{(i)}
    =
    (F_{t,t-1}^{l})^{-1}
    \left(\Xi_{t-1}^{(i)}\mid \widehat X_t^{(i)},\widehat\alpha^{(i)}\right).
    \tag{B2}
\]
This is Zhao--Cui Eq. (24).

Starting at final time \(T\), draw
\[
    (\widehat X_T,\widehat\alpha,\widehat X_{T-1})
    \sim
    \widehat p_T(x_T,\alpha,x_{T-1}\mid y_{1:T}),
\]
then move backward using (B2).  The resulting approximate path density factors
as Zhao--Cui Eq. (25):
\[
    \widehat p(\alpha,x_{0:T}\mid y_{1:T})
    =
    \widehat p_T(x_T,\alpha,x_{T-1}\mid y_{1:T})
    \prod_{t=1}^{T-1}
    \widehat p_t(x_{t-1}\mid x_t,\alpha,y_{1:t}).
    \tag{B3}
\]
The exact smoothing posterior has the same Markov factorization form with exact
conditionals:
\[
    p(\alpha,x_{0:T}\mid y_{1:T})
    =
    p(x_T,\alpha,x_{T-1}\mid y_{1:T})
    \prod_{t=1}^{T-1}
    p(x_{t-1}\mid x_t,\alpha,y_{1:t}).
    \tag{B4}
\]
The approximation can therefore be corrected by an importance ratio.  The
unnormalized backward weight, Zhao--Cui Eq. (26), is
\[
    \omega_b(\alpha,x_{0:T})
    =
    \frac{
    p(\alpha)p_\alpha(x_0)
    \prod_{t=1}^{T}
    f_\alpha(x_t\mid x_{t-1})g_\alpha(y_t\mid x_t)}
    {
    \widehat p_T(x_T,\alpha,x_{T-1}\mid y_{1:T})
    \prod_{t=1}^{T-1}
    \widehat p_t(x_{t-1}\mid x_t,\alpha,y_{1:t})}.
    \tag{B5}
\]
The numerator is the exact unnormalized path posterior.  The denominator is the
proposal density generated by the backward TT maps.

Algorithm 4 normalizes the backward weights in the same way:
\[
    W^{(i)}
    =
    \frac{\omega_b(\widehat\alpha^{(i)},\widehat x_{0:T}^{(i)})}
    {\sum_{j=1}^{N}\omega_b(\widehat\alpha^{(j)},\widehat x_{0:T}^{(j)})}.
    \tag{B6}
\]
The paper emphasizes a Markov identity used by the backward recursion:
\[
    p(x_{t-1}\mid \alpha,x_t,y_{1:t})
    =
    p(x_{t-1}\mid \alpha,x_t,y_{1:T}),
    \qquad t<T.
    \tag{B7}
\]
The reason is conditional independence.  Once \((\alpha,x_t)\) are known, the
future observations \(y_{t+1:T}\) depend on the past state \(x_{t-1}\) only
through \(x_t\).  This is why a conditional object built at time \(t\) can be
used inside a final-time smoothing path.

Zhao--Cui also give the computational reason for the variable order
\((x_t,\alpha,x_{t-1})\).  Marginalizing a squared TT from an end costs
\[
    O(pR^3)\quad\text{per coordinate},
    \tag{B8}
\]
because it updates one boundary mass matrix.  Marginalizing a coordinate in the
middle can cost
\[
    O(pR^6)\quad\text{per coordinate},
    \tag{B9}
\]
because left and right rank couplings must be paired.  The order
\((x_t,\alpha,x_{t-1})\) lets Algorithm 2 integrate \(x_{t-1}\) from the right
end and lets the smoothing map reuse the same direction.  The upper map for
particle filtering goes in the opposite direction and therefore incurs an
additional \(O(mpR^3)\)-type preprocessing cost.

Mini implementation contract.  Inputs: final-time approximation, lower
conditional maps for all times, model density evaluators, and reference uniform
samples.  Outputs: path samples and normalized smoothing weights.  Persisted
state: map sequence, samples, weights, and ESS.  Operation: final joint sample,
then backward inverse-map sampling.  Diagnostics: backward proposal support,
weight variance, root residuals, and mismatch between source-time conditionals
and final-time path use.

\section{Error Propagation: What It Proves And What It Does Not Prove}

Section 4 of Zhao--Cui explains how local approximation errors propagate.  The
key identity is the triangle inequality.  Let \(\pi_t\) be the exact
unnormalized joint density of \((x_t,\alpha,x_{t-1})\) at time \(t\), \(q_t\)
the propagated density using the previous approximation, and
\(\widehat\pi_t\) the new TT approximation.  Then
\[
    D(\widehat\pi_t,\pi_t)
    \le
    D(\widehat\pi_t,q_t)+D(q_t,\pi_t).
    \tag{E1}
\]
The first term is the new fit error.  The second is inherited from the previous
time step.  Under bounded likelihood or bounded prediction-likelihood
conditions, the inherited error is multiplied by a finite constant.

For squared TT, Zhao--Cui use an \(L^2\) error on square roots because it
controls Hellinger distance after normalization.  The practical message is
precise but modest: if every local square-root approximation is accurate enough
and the model does not amplify errors too severely, the sequence of approximate
densities can remain controlled.  This is not a theorem that ranks will stay
small, that the companion implementation is globally differentiable, or that a
particular high-dimensional model will be accurate without diagnostics.

\section{Preconditioning}

\paragraph{Source units S5.1-Eq30--Eq33.}
These source units answer: how can a hard target density be transformed into a
less concentrated residual before fitting a TT?

The direct target \(q_t(x_t,\alpha,x_{t-1})\) can be sharply concentrated.  A TT
may then need high rank.  Zhao--Cui Section 5 introduces a change of variables
that makes the function closer to a product reference density before fitting.

Let
\[
    r=(x_t,\alpha,x_{t-1}),\qquad u=T_t(r),
\]
and let \(\eta(u)=\eta_1(u_1)\cdots\eta_D(u_D)\) be a product reference density.
Choose a bridging density \(\rho_t(r)\) that is easier to approximate than
\(q_t(r)\), and construct \(T_t\) so that the pushforward of \(\rho_t\) is close
to \(\eta\):
\[
    (T_t)_\#\rho_t(u)
    =
    \rho_t(T_t^{-1}(u))
    \left|\det\nabla T_t^{-1}(u)\right|
    \propto
    \eta(u).
    \tag{P1}
\]
This is Zhao--Cui Eq. (30).

The pushforward of \(q_t\) is
\[
    (T_t)_\#q_t(u)
    =
    q_t(T_t^{-1}(u))
    \left|\det\nabla T_t^{-1}(u)\right|.
\]
Using (P1), the Jacobian factor is proportional to
\(\eta(u)/\rho_t(T_t^{-1}(u))\), so the function to fit is
\[
    q^\sharp_t(u)
    =
    \frac{q_t(T_t^{-1}(u))}
    {\rho_t(T_t^{-1}(u))}
    \eta(u).
    \tag{P2}
\]
This is Zhao--Cui Eq. (31).  If \(\rho_t\) captures the main shape of \(q_t\),
then \(q_t/\rho_t\) is flatter than \(q_t\), and TT ranks may be smaller.

In density language, the preconditioner should move the hard geometry into the
map \(T_t\).  The residual \(q_t^\sharp\) then asks the TT to learn only what
the bridge missed.  The implementation failure mode is division by a bridge
that is too small: \(\rho_t(T_t^{-1}(u))\) or \(\widehat\rho_t(T_t^{-1}(u))\)
must not vanish on target-relevant points.

Fit a square-root TT in \(u\)-space:
\[
    \sqrt{q^\sharp_t(u)}\approx \phi^\sharp_t(u),
\]
and define
\[
    \widehat\nu^\sharp_t(u)
    =
    \phi^\sharp_t(u)^2+\tau^\sharp_t\eta(u).
    \tag{P3}
\]
This is Zhao--Cui Eq. (32).  Pulling the approximation back to physical
coordinates gives
\[
    \widehat\pi_t(r)
    =
    \frac{\widehat\nu^\sharp_t(T_t(r))}
    {\eta(T_t(r))}
    \rho_t(r).
    \tag{P4}
\]
This is Zhao--Cui Eq. (33).  The normalizer is unchanged by the change of
variables:
\[
    \int \widehat\pi_t(r)\dd r
    =
    \int \widehat\nu^\sharp_t(u)\dd u.
    \tag{P5}
\]

Implementation contract for S5.1-Eq30--Eq33.  Inputs are \(q_t\), bridge
\(\rho_t\) or \(\widehat\rho_t\), map \(T_t\), reference density \(\eta\), and
residual TT \(\widehat\nu_t^\sharp\).  Outputs are a physical-coordinate
approximate density evaluator \(\widehat\pi_t\), a normalizer, and later
marginal/conditional map objects.  Persisted state: map evaluators, inverse map
evaluators, log-Jacobians or density ratios, bridge TT, residual TT, and
normalizer.  Diagnostics: bridge support ratio, map invertibility,
normalizer equality between physical and reference coordinates, and residual
rank.

\paragraph{Source unit S5.2-Gaussian-bridge.}
This source unit answers: what is the simplest concrete preconditioner?

A simple bridge is Gaussian.  Estimate a mean and covariance for \(r\), set
\(\rho_t=\mathcal N(\mu_t,\Sigma_t)\), and use an affine Gaussian whitening map.
If
\[
    \Sigma_t=L_tL_t^\top
    \tag{P6a}
\]
is a Cholesky factorization with \(L_t\) lower triangular, then the lower
linear KR whitening map is
\[
    T_t^{\ell}(r)
    =
    L_t^{-1}(r-\mu_t),
    \tag{P6b}
\]
and \(T_t^{\ell}(R)\sim N(0,I)\) when \(R\sim N(\mu_t,\Sigma_t)\).  The
reference density \(\eta\) is then the standard Gaussian product density.  The
implementation stores \(\mu_t\), \(L_t\), triangular solves for \(L_t^{-1}\),
and the log determinant
\[
    \log|\det\nabla T_t^{\ell}|
    =
    -\sum_j\log L_t[j,j].
    \tag{P6c}
\]
An upper triangular linear map is obtained by permuting or reversing the
coordinate order before applying the same whitening idea.

Implementation contract for S5.2.  Inputs: particles or quadrature samples used
to estimate \(\mu_t,\Sigma_t\).  Outputs: triangular factor \(L_t\), whitening
map \(T_t^\ell\), inverse map \(r=\mu_t+L_tu\), and log determinant.  Persisted
state: \(\mu_t,L_t\) and any covariance regularization.  Diagnostics:
positive-definiteness of \(\Sigma_t\), conditioning of \(L_t\), and whether the
Gaussian bridge assigns negligible density to target-relevant samples.

\paragraph{Source unit S5.3-Eq34--Eq35.}
This source unit answers: how can the bridge retain nonlinear structure from
the previous filter, transition, and likelihood?

A nonlinear bridge uses powers of the three factors in \(q_t\):
\[
    \rho_t(r)
    =
    \widehat p_{t-1}(x_{t-1},\alpha)^{\beta_\pi}
    f_\alpha(x_t\mid x_{t-1})^{\beta_f}
    g_\alpha(y_t\mid x_t)^{\beta_g},
    \qquad
    \beta_\pi,\beta_f,\beta_g\in[0,1].
    \tag{P6}
\]
This is Zhao--Cui Eq. (34).  A squared-TT approximation of this bridge is
\[
    \widehat\rho_t(r)
    =
    \phi_{\rho,t}(r)^2+\tau_{\rho,t}\lambda(r).
    \tag{P7}
\]
This is Zhao--Cui Eq. (35).  The same algebra as (P4) applies with
\(\rho_t\) replaced by \(\widehat\rho_t\).

Implementation contract for S5.3.  Inputs: exponents
\(\beta_\pi,\beta_f,\beta_g\), previous filter evaluator, transition evaluator,
likelihood evaluator, and bridge TT controls.  Outputs: \(\widehat\rho_t\) and
its KR map.  Persisted state: exponents, bridge TT cores, defensive mass, and
bridge mass contractions.  Diagnostics: bridge fit residual, bridge rank,
support coverage, and whether the powered factors are finite at fitting points.

\subsection{The Five Densities In The Preconditioned Step}

\paragraph{Source unit S5.3-density-chain.}
This source unit expands the compact Section 5 notation into the five densities
an implementer must distinguish.

The preconditioned step is easy to lose because the paper moves between
physical coordinates \(r=(x_t,\theta,x_{t-1})\), reference coordinates
\(u=(u_t,u_\theta,u_{t-1})\), an exact bridge, an approximate bridge, and a
residual density.  The complete chain is as follows.

First, there is the exact target we ultimately want to approximate:
\[
    q_t(r)
    =
    \widehat\pi_{t-1}(x_{t-1},\theta)
    f_\theta(x_t\mid x_{t-1})
    g_\theta(y_t\mid x_t).
    \tag{P7a}
\]
Second, there is an exact bridge \(\rho_t(r)\).  It is not the final answer.
It is a simpler density chosen to resemble the main shape of \(q_t\).  In the
tempered construction it is the powered product in (P6).  In the Gaussian
construction it is \(N(\mu_t,\Sigma_t)\).

Third, because the exact bridge may still be nonseparable, Algorithm 5(b.1)
approximates the bridge itself by a squared TT:
\[
    \sqrt{\rho_t(r)}\approx \phi_{\rho,t}(r),
    \qquad
    \widehat\rho_t(r)
    =
    \phi_{\rho,t}(r)^2+\tau_{\rho,t}\lambda(r).
    \tag{P7b}
\]
This is the first line of Algorithm 5 in mathematical form.  The implementation
stores the bridge TT cores, the defensive mass \(\tau_{\rho,t}\), the reference
density \(\lambda\), the bridge normalizer, and the mass contractions needed to
construct conditional maps from \(\widehat\rho_t\).

Fourth, build the preconditioning map from the bridge approximation.  Let
\[
    T_t:\ r\mapsto u
    \tag{P7c}
\]
be a KR-based map satisfying
\[
    (T_t)_\#\widehat\rho_t(u)\propto \eta(u).
    \tag{P7d}
\]
Using \(\widehat\rho_t\), not the unavailable exact bridge, the residual target
to approximate in reference coordinates is
\[
    q_t^\sharp(u)
    =
    \frac{q_t(T_t^{-1}(u))}
    {\widehat\rho_t(T_t^{-1}(u))}
    \eta(u).
    \tag{P7e}
\]
This formula answers the question ``what remains after the bridge has removed
the main shape?''  If the bridge is good, the ratio \(q_t/\widehat\rho_t\) is
flatter than \(q_t\), so the residual may need smaller TT ranks.

Fifth, approximate the square root of the residual:
\[
    \sqrt{q_t^\sharp(u)}\approx \phi_t^\sharp(u),
    \qquad
    \widehat\nu_t^\sharp(u)
    =
    \phi_t^\sharp(u)^2+\tau_t^\sharp\eta(u).
    \tag{P7f}
\]
Finally pull this residual approximation back to physical coordinates:
\[
    \widehat\pi_t(r)
    =
    \frac{\widehat\nu_t^\sharp(T_t(r))}
    {\eta(T_t(r))}
    \widehat\rho_t(r).
    \tag{P7g}
\]
To check the algebra, substitute \(u=T_t(r)\).  Since
\((T_t)_\#\widehat\rho_t\propto\eta\), the factor
\(\widehat\rho_t(r)/\eta(T_t(r))\) restores the physical-coordinate density
scale.  The normalizer is computed in reference space:
\[
    \int \widehat\pi_t(r)\dd r
    =
    \int \widehat\nu_t^\sharp(u)\dd u.
    \tag{P7h}
\]
Algorithm 5 stores exactly the objects in this chain:
\[
    \widehat\rho_t,\quad T_t,\quad q_t^\sharp\text{ evaluator},\quad
    \widehat\nu_t^\sharp,\quad \widehat\pi_t\text{ evaluator},\quad
    \text{marginal and conditional maps}.
    \tag{P7i}
\]
This is the preconditioning story in one line of thought: approximate the
bridge, use the bridge to flatten the target, approximate the flattened
residual, then pull the residual back.

Failure checks for the five-density chain.  Verify all of the following:
\[
    q_t(r)\ge0,\qquad
    \widehat\rho_t(r)>0\ \text{on fitted target points},\qquad
    q_t^\sharp(u)\ge0,
    \tag{P7j}
\]
\[
    0<\int\widehat\nu_t^\sharp(u)\dd u<\infty,\qquad
    \left|
    \int\widehat\pi_t(r)\dd r-\int\widehat\nu_t^\sharp(u)\dd u
    \right|
    \le \varepsilon_{\rm norm}.
    \tag{P7k}
\]
These are not cosmetic checks.  They assess whether the density chain still
defines the same scalar after moving between coordinate systems.

\paragraph{Source unit S5.4-Algorithm5(b.1).}
Algorithm 5(b.1) answers: what approximation is built before the residual
target is fitted?

The bridge approximation is \(\widehat\rho_t\), already defined in (P7b).  Its
mathematical input-output contract is
\[
    (r\mapsto \rho_t(r),\lambda,\tau_{\rho,t},\text{basis},\text{ranks})
    \mapsto
    (\widehat\rho_t,\text{bridge mass contractions}).
    \tag{P7l}
\]
Failure checks: bridge fit residual, positivity of \(\widehat\rho_t\), and
nonzero bridge values on target fitting points.

\paragraph{Source unit S5.4-Algorithm5(b.2).}
Algorithm 5(b.2) answers: how does the bridge approximation become a lower
preconditioning map and a bridge marginal?

Build a KR map \(R_t\) from \(\widehat\rho_t\) to a uniform variable:
\[
    (R_t)_\#\widehat\rho_t(\xi_t,\xi_\theta,\xi_{t-1})
    \propto
    \operatorname{Unif}(\xi_t,\xi_\theta,\xi_{t-1}).
    \tag{P8}
\]
If \(D\) is a diagonal coordinate map satisfying
\[
    D_\#\eta(\xi_t,\xi_\theta,\xi_{t-1})
    =
    \operatorname{Unif}(\xi_t,\xi_\theta,\xi_{t-1}),
    \tag{P9}
\]
then
\[
    T_t=D^{-1}\circ R_t
    \tag{P10}
\]
pushes the bridge approximation to the product reference \(\eta\):
\[
    (T_t)_\#\widehat\rho_t(u_t,u_\theta,u_{t-1})
    \propto
    \eta(u_t,u_\theta,u_{t-1}).
    \tag{P11}
\]
For the lower-triangular version, write
\[
    T_t^{\ell}(x_t,\theta,x_{t-1})
    =
    \bigl(
    T_{t,t}^{\ell}(x_t),
    T_{t,\theta}^{\ell}(\theta\mid x_t),
    T_{t,t-1}^{\ell}(x_{t-1}\mid x_t,\theta)
    \bigr).
    \tag{P11a}
\]
The same bridge mass contractions also provide the bridge marginal
\[
    \widehat\rho_t(x_t,\theta)
    =
    \int\widehat\rho_t(x_t,\theta,x_{t-1})\dd x_{t-1}.
    \tag{P11b}
\]
Mathematical input-output contract:
\[
    \widehat\rho_t
    \mapsto
    (T_t^\ell,(T_t^\ell)^{-1},\widehat\rho_t(x_t,\theta),
    T_{t,t-1}^\ell).
    \tag{P11c}
\]
Failure checks: monotone conditional CDFs, inverse-map residuals, and agreement
between the bridge marginal and direct bridge contraction.

\paragraph{Source unit S5.4-Algorithm5(b.3).}
Algorithm 5(b.3) answers: after the bridge map is available, what residual
density is fitted in reference coordinates?

With this nonlinear preconditioner, the residual target becomes
\[
    q^\sharp_t(u)
    =
    \frac{q_t(T_t^{-1}(u))}
    {\widehat\rho_t(T_t^{-1}(u))}
    \eta(u),
    \tag{P12}
\]
and the pulled-back unnormalized posterior approximation is
\[
    \widehat\pi_t(r)
    =
    \frac{\widehat\nu^\sharp_t(T_t(r))}
    {\eta(T_t(r))}
    \widehat\rho_t(r).
    \tag{P13}
\]
Mathematical input-output contract:
\[
    (q_t,T_t^{-1},\widehat\rho_t,\eta,\text{basis},\text{ranks})
    \mapsto
    (\widehat\nu_t^\sharp,S_t^\ell,\text{residual mass contractions}).
    \tag{P13a}
\]
Failure checks: finite ratio \(q_t/\widehat\rho_t\), residual rank growth, and
normalizer \(\int\widehat\nu_t^\sharp(u)\dd u\in(0,\infty)\).

\paragraph{Source unit S5.4-Algorithm5(c.1).}
Algorithm 5(c.1) answers: what residual marginal is needed before returning to
physical coordinates?

Use
the lower triangular decompositions
\[
    S_t^{\ell}(u_t,u_\theta,u_{t-1})
    =
    \begin{bmatrix}
    S_{t,t}^{\ell}(u_t)\\
    S_{t,\theta}^{\ell}(u_\theta\mid u_t)\\
    S_{t,t-1}^{\ell}(u_{t-1}\mid u_t,u_\theta)
    \end{bmatrix},
    \qquad
    T_t^{\ell}(x_t,\theta,x_{t-1})
    =
    \begin{bmatrix}
    T_{t,t}^{\ell}(x_t)\\
    T_{t,\theta}^{\ell}(\theta\mid x_t)\\
    T_{t,t-1}^{\ell}(x_{t-1}\mid x_t,\theta)
    \end{bmatrix}.
    \tag{P14}
\]
Step (c.1) integrates the last block \((u_{t-1})\) of
\(\widehat\nu_t^\sharp\) to obtain
\[
    \widehat\nu_t^\sharp(u_t,u_\theta)
    =
    \int\widehat\nu_t^\sharp(u_t,u_\theta,u_{t-1})\dd u_{t-1}.
    \tag{P15}
\]
This is an endpoint squared-TT marginalization in \(u\)-space, so it uses the
same mass-matrix algebra as Proposition 2.  It also creates the residual lower
conditional map \(S_{t,t-1}^{\ell}(u_{t-1}\mid u_t,u_\theta)\).

\paragraph{Source unit S5.4-Algorithm5(c.2).}
Algorithm 5(c.2) answers: how does the residual marginal become the retained
physical-coordinate filter?

Using
\[
    (u_t,u_\theta)
    =
    \bigl(
    T_{t,t}^{\ell}(x_t),
    T_{t,\theta}^{\ell}(\theta\mid x_t)
    \bigr),
    \tag{P16}
\]
start from the pulled-back three-block density (P13):
\[
    \widehat\pi_t(x_t,\theta,x_{t-1})
    =
    \frac{
    \widehat\nu_t^\sharp(T_t^\ell(x_t,\theta,x_{t-1}))}
    {\eta(T_t^\ell(x_t,\theta,x_{t-1}))}
    \widehat\rho_t(x_t,\theta,x_{t-1}).
    \tag{P16a}
\]
Now change variables only in the old-state conditional block.  For fixed
\((x_t,\theta)\), the lower bridge map sends
\[
    x_{t-1}
    \mapsto
    u_{t-1}
    =
    T_{t,t-1}^{\ell}(x_{t-1}\mid x_t,\theta).
    \tag{P16b}
\]
Because \(T_t^\ell\) was built from the bridge approximation, the bridge
density decomposes locally.  Write
\[
    u_t=T_{t,t}^{\ell}(x_t),\qquad
    u_\theta=T_{t,\theta}^{\ell}(\theta\mid x_t).
    \tag{P16b.1}
\]
The conditional part of the triangular pushforward identity says that, for
fixed \((x_t,\theta)\),
\[
    \widehat\rho_t(x_{t-1}\mid x_t,\theta)\dd x_{t-1}
    =
    \eta(u_{t-1}\mid u_t,u_\theta)\dd u_{t-1}.
    \tag{P16b.2}
\]
Since the reference density factorizes in the transformed coordinates,
\[
    \eta(u_t,u_\theta,u_{t-1})
    =
    \eta(u_{t-1}\mid u_t,u_\theta)\eta(u_t,u_\theta),
    \tag{P16b.3}
\]
the conditional density in (P16b.2) can be written as
\[
    \eta(u_{t-1}\mid u_t,u_\theta)
    =
    \frac{\eta(T_t^\ell(x_t,\theta,x_{t-1}))}
    {\eta(T_{t,t}^{\ell}(x_t),T_{t,\theta}^{\ell}(\theta\mid x_t))}.
    \tag{P16b.4}
\]
Multiplying (P16b.2) by the bridge marginal
\(\widehat\rho_t(x_t,\theta)\) gives the local bridge decomposition
\[
    \widehat\rho_t(x_t,\theta,x_{t-1})\dd x_{t-1}
    =
    \frac{
    \widehat\rho_t(x_t,\theta)}
    {\eta(T_{t,t}^{\ell}(x_t),T_{t,\theta}^{\ell}(\theta\mid x_t))}
    \eta(T_t^\ell(x_t,\theta,x_{t-1}))\dd u_{t-1},
    \tag{P16c}
\]
where the Jacobian of the conditional map is included in the differential
\(\dd u_{t-1}\).  This is the missing middle step in plain words: the
three-block bridge density equals its two-block marginal times the conditional
bridge density, and the conditional bridge density becomes the conditional
reference density after the triangular map.  Substituting (P16c) into (P16a)
and integrating over
\(x_{t-1}\) gives
\[
\begin{aligned}
 \widehat\pi_t(x_t,\theta\mid y_{1:t})
 &=
 \int \widehat\pi_t(x_t,\theta,x_{t-1})\dd x_{t-1}\\
 &=
 \frac{\widehat\rho_t(x_t,\theta)}
 {\eta(T_{t,t}^{\ell}(x_t),T_{t,\theta}^{\ell}(\theta\mid x_t))}
 \int \widehat\nu_t^\sharp(u_t,u_\theta,u_{t-1})\dd u_{t-1}.
\end{aligned}
    \tag{P16d}
\]
Using (P15) with
\((u_t,u_\theta)=
(T_{t,t}^{\ell}(x_t),T_{t,\theta}^{\ell}(\theta\mid x_t))\), the retained
marginal density for the next filtering step is
\[
    \widehat\pi_t(x_t,\theta\mid y_{1:t})
    =
    \frac{
    \widehat\nu_t^\sharp\!\left(
    T_{t,t}^{\ell}(x_t),
    T_{t,\theta}^{\ell}(\theta\mid x_t)
    \right)}
    {
    \eta\!\left(
    T_{t,t}^{\ell}(x_t),
    T_{t,\theta}^{\ell}(\theta\mid x_t)
    \right)}
    \widehat\rho_t(x_t,\theta).
    \tag{P17}
\]
This is the source formula in Algorithm 5(c.2) written without relying on
paper shorthand.  The implementation stores the two maps, their marginal
evaluators, and a retained evaluator for \(\widehat\pi_t(x_t,\theta)\).

As a byproduct, the lower conditional map for path estimation is the last block
of the composite map:
\[
    F_{t,t-1}^{\ell}(x_{t-1}\mid x_t,\theta)
    =
    S_{t,t-1}^{\ell}
    \!\left(
    T_{t,t-1}^{\ell}(x_{t-1}\mid x_t,\theta)
    \,\middle|\,
    T_{t,t}^{\ell}(x_t),
    T_{t,\theta}^{\ell}(\theta\mid x_t)
    \right).
    \tag{P18}
\]
This identity matters because preconditioning does not eliminate smoothing; it
changes which conditional maps must be composed and stored.

\begin{center}
\fbox{\begin{minipage}{0.94\linewidth}
\textbf{Boxed Algorithm 5 dataflow.}
\[
\begin{array}{llll}
\text{Step} & \text{Evaluator signature} & \text{Persisted objects} & \text{Main diagnostic}\\
\hline
\text{b.1} &
\rho_t(r)\mapsto\widehat\rho_t(r) &
\text{bridge TT, }\tau_{\rho,t},\lambda &
\text{bridge residual/rank}\\
\text{b.2} &
\widehat\rho_t\mapsto T_t^\ell,(T_t^\ell)^{-1} &
\text{bridge maps, bridge marginal} &
\text{CDF monotonicity}\\
\text{b.3} &
(q_t,T_t^{-1},\widehat\rho_t,\eta)\mapsto\widehat\nu_t^\sharp &
\text{residual TT, }S_t^\ell &
\text{ratio finite}\\
\text{c.1} &
\widehat\nu_t^\sharp(u_t,u_\theta,u_{t-1})\mapsto
\widehat\nu_t^\sharp(u_t,u_\theta) &
\text{residual marginal, }S_{t,t-1}^\ell &
\text{mass contraction}\\
\text{c.2} &
(\widehat\nu_t^\sharp,T_t^\ell,\widehat\rho_t)\mapsto
\widehat\pi_t(x_t,\theta) &
\text{retained filter evaluator} &
\text{normalizer invariance}
\end{array}
\]
\end{minipage}}
\end{center}

\section{End of Zhao--Cui Annotation and Start of BayesFilter Fixed-Branch Extension}

Up to this point the note has annotated Zhao--Cui Sections 1--3 and 5 in
source order.  The equations above explain the paper's adaptive sequential
algorithm: tensor-train approximation, squared-TT nonnegativity, conditional KR
maps, particle correction, smoothing, and preconditioning.  The next sections
are a BayesFilter extension.  They freeze approximation choices so that an
analytical derivative can differentiate the same scalar computed by the forward
pass.  Zhao--Cui's adaptive rank, pivot, fitting, preconditioning, and branch
choices should not be read as globally differentiable.

\subsection*{Reconciliation: Paper Objects Versus Fixed-Branch Objects}

The paper's main sequential-learning density carries a learned static
parameter as a random coordinate:
\[
    r_t^{\rm ZC}=(x_t,\theta,x_{t-1}).
    \tag{R1}
\]
The fixed-branch likelihood lane conditions on an external parameter value
\(\beta\) and differentiates with respect to that external argument:
\[
    r_t^{\rm FB}=(x_t,x_{t-1};\beta).
    \tag{R2}
\]
The semicolon is deliberate: \(\beta\) is not integrated out by the TT density.
It is the input at which the approximate likelihood is evaluated.

The exact mapping of objects is:
\[
\begin{gathered}
\widehat\pi_{t-1}(x_{t-1},\theta)
\mapsto
\widehat p_{t-1}(x_{t-1};\beta),\\
f(x_t\mid x_{t-1},\theta)\mapsto f_\beta(x_t\mid x_{t-1}),
\qquad
g(y_t\mid x_t,\theta)\mapsto g_\beta(y_t\mid x_t),\\
\phi_t(x_t,\theta,x_{t-1})
\mapsto
\phi_t(x_t,x_{t-1};\beta),\\
\tau_t\lambda(x_t)\lambda(\theta)\lambda(x_{t-1})
\mapsto
\tau_t\lambda_t(x_t,x_{t-1}),\\
\widehat Z_t
=
\int\widehat q_t(x_t,\theta,x_{t-1})
\mapsto
\widehat Z_t(\beta)
=
\int\widehat q_t(x_t,x_{t-1};\beta),\\
\text{adaptive ranks/pivots/maps}
\mapsto
\text{fixed ranks/points/maps/sweeps}.
\end{gathered}
\tag{R3}
\]
The first line keeps the retained-filter role but drops the integrated
\(\theta\) coordinate.  The second line keeps the model evaluators, evaluated at
the external \(\beta\).  The third and fourth lines keep the squared-TT and
defensive-density forms in a lower-dimensional state space.  The fifth line
keeps the evidence-increment role.  The last line is the essential change:
adaptive choices are replaced by frozen branch choices before differentiation.
Thus the fixed-branch section does not differentiate the adaptive Zhao--Cui
algorithm.  It differentiates a lower-dimensional scalar constructed by applying
the same squared-TT filtering algebra after freezing all discrete approximation
choices.  If one wants a fixed-branch derivative while also learning a static
parameter as a random coordinate, one can append that coordinate to the state;
the external derivative variable must still be distinguished from the
integrated coordinate.

\section{Implementable Fixed-Branch Objects And Data Structures}

A reader implementing a minimal method must store more than an equation.  At
time \(t\), a squared-TT filter object contains:
\[
    \mathcal B_t
    =
    \{
    \Omega_{t,k},\Psi_{t,k},\psi_{t,k,\ell},C_{t,k},
    R_{t,k},B_{t,k},\lambda_{t,k},\tau_t,c_t,
    \text{mass contractions},\text{map data}
    \}.
\]
Here \(\Omega_{t,k}\) is the coordinate domain, \(\Psi_{t,k}\) maps reference
coordinates to physical coordinates, \(\psi_{t,k,\ell}\) are basis functions,
\(C_{t,k}\) are TT cores, \(R_{t,k}\) are ranks, \(B_{t,k}\) are mass matrices,
\(\lambda_{t,k}\) are reference densities, \(\tau_t\) is the defensive mass, and
\(c_t\) is any log-scaling shift used during fitting.

For the adaptive Zhao--Cui algorithm these objects may change with the data,
with samples, and with approximation decisions.  For a differentiable
same-scalar target, the realized branch must be frozen before differentiation.

\subsection{What A Core Object Looks Like}

For coordinate \(k\), store
\[
    C_k\in\R^{R_{k-1}\times p_k\times R_k}.
\]
Given a scalar coordinate value \(r_k\), first compute the basis vector
\[
    b_k(r_k)=
    \bigl(\psi_{k,0}(r_k),\ldots,\psi_{k,p_k-1}(r_k)\bigr)^\top.
\]
Then form the matrix
\[
    H_k(r_k)[a,b]
    =
    \sum_{\ell=0}^{p_k-1}C_k[a,\ell,b]\,b_k(r_k)[\ell].
    \tag{D1}
\]
Evaluation of the TT is the loop
\[
    v_0=1,\qquad
    v_k=v_{k-1}H_k(r_k),\quad k=1,\ldots,D,
    \qquad
    \widehat h(r)=v_D.
    \tag{D2}
\]
The square-density object additionally stores mass contractions
\[
    M_{>k}
    =
    \int H_{k+1}(r_{k+1})\cdots H_D(r_D)
    H_D(r_D)^\top\cdots H_{k+1}(r_{k+1})^\top
    \dd r_{k+1:D}.
    \tag{D3}
\]
These matrices are enough to evaluate marginals from the left.  Analogous
left-side mass contractions evaluate reverse-order marginals.

\subsection{What Must Be Saved For The Next Time Step}

After time \(t\), the next step needs to evaluate
\[
    \widehat p_t(x_t;\beta)
    \quad\text{and}\quad
    \partial_{\beta_i}\widehat p_t(x_t;\beta)
\]
at future fitting points.  It is not enough to store a cloud of samples or only
\(\widehat Z_t\).  A fixed-branch derivative implementation stores:
\[
    \left\{
    C_{t,k},\dot C_{t,k}^{(i)},M_{t,>k},\dot M_{t,>k}^{(i)},
    \widehat Z_t,\dot{\widehat Z}_t^{(i)},
    \Psi_t,\lambda_t,\tau_t
    \right\}.
    \tag{D4}
\]
The derivative of the normalized marginal has the quotient form
\[
    \partial_{\beta_i}\widehat p_t(x_t;\beta)
    =
    \frac{\partial_{\beta_i}a_t(x_t;\beta)}{\widehat Z_t(\beta)}
    -
    \frac{a_t(x_t;\beta)\partial_{\beta_i}\widehat Z_t(\beta)}
    {\widehat Z_t(\beta)^2},
    \tag{D5}
\]
where
\[
    a_t(x_t;\beta)
    =
    \int \widehat q_t(x_t,x_{t-1};\beta)\dd x_{t-1}
    \tag{D6}
\]
is the unnormalized retained numerator.  Equation (D5) is the carried-filter
derivative needed in (G2) at the next time step.

\begin{lemma}[Carried filter object feeds the next target]
Fix the branch at time \(t\).  Suppose the stored object contains
\[
    \mathcal F_t
    =
    \{a_t,\dot a_t^{(i)},\widehat Z_t,\dot{\widehat Z}_t^{(i)},
      \Psi_t,\lambda_t,\tau_t\}.
\]
Define \(\widehat p_t\) and \(\dot{\widehat p}_t^{(i)}\) by (D5).  Then the
next transformed target and its derivative at any next-step fitting point are
computable from \(\mathcal F_t\), the model density evaluators, and their
derivatives.
\end{lemma}

\begin{proof}
At time \(t+1\), the fixed-branch target has the form
\[
    \widetilde q_{t+1}(z;\beta)
    =
    \widehat p_t(x_t(z);\beta)
    f_\beta(x_{t+1}(z)\mid x_t(z))
    g_\beta(y_{t+1}\mid x_{t+1}(z))
    J_{t+1}(z).
\]
The stored object evaluates \(\widehat p_t\).  The model supplies \(f_\beta\)
and \(g_\beta\), and the branch supplies \(J_{t+1}\).  Differentiating gives
the product-rule expression (G2) with \(t\) replaced by \(t+1\).  The first
term uses \(\dot{\widehat p}_t^{(i)}\), which is the quotient derivative in
(D5).  No old samples or external source-paper objects are needed.
\end{proof}

\section{One Concrete Branch, Fully Instantiated}

This section gives a single implementation lane.  It is intentionally narrower
than Zhao--Cui's adaptive code.  Its purpose is to remove ambiguity.  A coder
can implement this lane first, then replace pieces with adaptive or higher
performance choices later.

\subsection{Coordinate Order And Domain}

For the fixed-parameter likelihood case, use the two-block coordinate order
\[
    r_t=(x_t,x_{t-1})\in\R^2
\]
for a scalar state.  For vector states, concatenate coordinates as
\[
    r_t=(x_{t,1},\ldots,x_{t,m_x},x_{t-1,1},\ldots,x_{t-1,m_x}).
\]
Choose a fixed interval for each coordinate,
\[
    r_{t,k}\in[\ell_{t,k},u_{t,k}],
\]
and map from reference coordinate \(z_k\in[-1,1]\) by
\[
    r_{t,k}=a_{t,k}+h_{t,k}z_k,\qquad
    a_{t,k}=\frac{u_{t,k}+\ell_{t,k}}2,\qquad
    h_{t,k}=\frac{u_{t,k}-\ell_{t,k}}2.
    \tag{C1}
\]
The Jacobian is constant:
\[
    J_t=\prod_{k=1}^{D}h_{t,k}.
    \tag{C2}
\]
The branch stores \(\ell_{t,k},u_{t,k},a_{t,k},h_{t,k}\).  In the same-scalar
derivative these are constants and \(\dot J_t=0\).

\subsection{Basis, Points, Weights, Ranks}

Use normalized Legendre basis with the same degree \(p\) in every coordinate:
\[
    b_k(z_k)=
    \left(
    \sqrt{\frac12}P_0(z_k),
    \sqrt{\frac32}P_1(z_k),
    \ldots,
    \sqrt{\frac{2p-1}{2}}P_{p-1}(z_k)
    \right)^\top.
    \tag{C3}
\]
Use fixed Halton-type fitting points.  Let \(p_k^\star\) be the \(k\)-th prime.
If
\[
    j=\sum_{m=0}^{M}d_m(p_k^\star)^m,\qquad
    d_m\in\{0,\ldots,p_k^\star-1\},
\]
define
\[
    \operatorname{radinv}_{p_k^\star}(j)
    =
    \sum_{m=0}^{M}d_m(p_k^\star)^{-(m+1)},
    \qquad
    z_k^{(j)}=2\operatorname{radinv}_{p_k^\star}(j)-1.
    \tag{C4}
\]
Use \(j=1,\ldots,N_{\rm fit}\), weights
\[
    W=\frac1{N_{\rm fit}}I,
    \tag{C5}
\]
and fixed ranks
\[
    R_0=R_D=1,\qquad R_1,\ldots,R_{D-1}\ \text{declared before fitting}.
    \tag{C6}
\]
No pivot, point, rank, or domain choice is allowed to change during a
same-scalar derivative check.

\subsection{Defensive Reference, Mass, And Shift}

Use the uniform reference on \([-1,1]^D\):
\[
    \lambda(z)=2^{-D}.
    \tag{C7}
\]
Choose a defensive floor \(\epsilon_\tau>0\) and set
\[
    \tau_t=2^D\epsilon_\tau.
    \tag{C8}
\]
Then \(\tau_t\lambda(z)=\epsilon_\tau\).  Choose the stabilizing shift at the
base parameter \(\beta_0\):
\[
    c_t=-\max_{1\le j\le N_{\rm fit}}
    \log \widetilde q_t(z^{(j)};\beta_0).
    \tag{C9}
\]
The fitted square-root target is
\[
    y_j(\beta)=\exp(c_t/2)\sqrt{\widetilde q_t(z^{(j)};\beta)}.
    \tag{C10}
\]
For same-branch finite differences, \(c_t\) is reused for
\(\beta_0+h\) and \(\beta_0-h\).  Recomputing \(c_t\) changes the scalar.

\subsection{Boxed Convention: Shifted Fit, Unshifted Density}

The concrete branch uses this convention everywhere downstream.
\[
\boxed{
\begin{aligned}
&\text{Fitted TT:}\qquad
  \phi_t(z;\beta)\approx e^{c_t/2}\sqrt{\widetilde q_t(z;\beta)},\\
&\text{Approximate joint density:}\qquad
  \widehat q_t(z;\beta)
  =
  e^{-c_t}\phi_t(z;\beta)^2+\tau_t\lambda_t(z),\\
&\text{Retained numerator:}\qquad
  a_t(z_t;\beta)
  =
  \int
  \left[e^{-c_t}\phi_t(z_t,z_{t-1};\beta)^2
        +\tau_t\lambda_t(z_t,z_{t-1})\right]\dd z_{t-1},\\
&\text{Normalizer:}\qquad
  \widehat Z_t(\beta)
  =
  \int \widehat q_t(z;\beta)\dd z
  =
  e^{-c_t}\int\phi_t(z;\beta)^2\dd z+\tau_t.
\end{aligned}}
\tag{C11}
\]
Every carried filter, proposition, derivative, and finite-difference check uses
(C11).  If an implementation instead absorbs \(e^{-c_t/2}\) into the TT cores,
then it must remove the explicit \(e^{-c_t}\) factors everywhere.  Mixing the
two conventions changes the scalar.

\subsection{Initialization And Sweep Order}

Initialize every core to zero except the constant channel:
\[
    C_k[1,0,1]=1,\qquad k=2,\ldots,D,
\]
and
\[
    C_1[1,0,1]=\bar y,\qquad
    \bar y=\frac1{N_{\rm fit}}\sum_{j=1}^{N_{\rm fit}}y_j.
\]
If a rank is larger than one, all non-first rank channels start at zero.  This
plain initialization gives a constant initial approximation and avoids a random
initialization branch.

One bidirectional sweep means update cores \(1,2,\ldots,D\), then
\(D,D-1,\ldots,1\).  Fix the number of bidirectional sweeps \(S\).  The core
update is the ridge solve (A1-2d).  The branch stores the sweep order, \(S\),
\(\rho\), and the final cores.

\subsection{Forward-Step Object Flow}

The fixed-branch forward step is a dependency chain among mathematical
objects, not executable code.  The stored-object flow is
\[
\begin{array}{rcl}
\text{reference points} &\longrightarrow& z^{(j)} \text{ from (C4)},\\
\text{physical points} &\longrightarrow& r^{(j)}=\Psi_t(z^{(j)}),\\
\text{target values} &\longrightarrow&
    \widetilde q_{t,j}
    =
    \widehat p_{t-1}(x_{t-1}^{(j)};\beta)
    f_\beta(x_t^{(j)}\mid x_{t-1}^{(j)})
    g_\beta(y_t\mid x_t^{(j)})J_t,\\
\text{square-root fit values} &\longrightarrow&
    y_j=\exp(c_t/2)\sqrt{\widetilde q_{t,j}},\\
\text{initial cores} &\longrightarrow&
    \text{constant-channel initialization},\\
\text{fixed sweep objects} &\longrightarrow&
    L_{j,k},R_{j,k},A_{j,k},N_k,d_k,C_k,\\
\text{mass objects} &\longrightarrow&
    M_{>k},R_t,\widehat Z_t,\\
\text{carried objects} &\longrightarrow&
    a_t,\widehat p_t.
\end{array}
\tag{C12}
\]
The fixed sweep environments are
\[
    L_{j,k}=H_1(z_1^{(j)})\cdots H_{k-1}(z_{k-1}^{(j)}),
\]
\[
    R_{j,k}=H_{k+1}(z_{k+1}^{(j)})\cdots H_D(z_D^{(j)}),
\]
with \(A_{j,k}\) defined by (A1-2b) and the core update defined by (A1-2d).
After fitting, square-mass contractions use (M5)--(M6):
\[
    R_t=\int\phi_t(z)^2\dd z,\qquad
    \widehat Z_t=e^{-c_t}R_t+\tau_t.
\]
The carried numerator is the old-state marginal
\[
    a_t(z_t;\beta)
    =
    \int
    \left[e^{-c_t}\phi_t(z_t,z_{t-1};\beta)^2
    +\tau_t\lambda_t(z_t,z_{t-1})\right]\dd z_{t-1},
\]
and the carried filter is \(a_t/\widehat Z_t\).

\subsection{Conditional-Map Inversion Protocol}

For a conditional coordinate \(r_k\), compute
\[
    F_k(s\mid r_{<k})
    =
    \int_{\ell_k}^{s}\widehat p(r_k'\mid r_{<k})\dd r_k'.
\]
To sample from the conditional, given \(u\in(0,1)\), solve
\[
    F_k(s\mid r_{<k})-u=0,\qquad s\in[\ell_k,u_k].
\]
Use bisection as the baseline because it only needs monotonicity.  Fix a
tolerance \(\varepsilon_{\rm root}\) and maximum iterations \(N_{\rm root}\).
The branch stores both.  The bisection update is:
\[
    m=\frac{a+b}{2},\qquad
    \text{if }F_k(m\mid r_{<k})<u\text{ set }a=m,\text{ else set }b=m.
\]
Stop when \(b-a\le\varepsilon_{\rm root}\) or the iteration count reaches
\(N_{\rm root}\).  Raise a failure diagnostic if the final residual
\(|F_k(m\mid r_{<k})-u|\) exceeds the declared tolerance.

\section{Fixed-Branch TT Filtering Recursion}

The fixed-branch BayesFilter variant chooses all discrete approximation objects
before differentiating.  It is narrower than the full Zhao--Cui adaptive code,
but it is the variant for which an analytical derivative can be stated without
changing the scalar under the derivative.

There is one important notational separation.  In the Zhao--Cui learning
algorithm, the unknown parameter \(\alpha\) may be carried as a random variable
inside the posterior density.  In an HMC-style likelihood calculation, the
parameter is instead an external argument that is varied by the sampler.  To
avoid using the same symbol for both roles, this section writes the external
differentiated parameter as
\[
    \beta\in\R^{m_\beta}.
\]
The fixed-branch likelihood variant conditions on a trial \(\beta\) and filters
only the latent state.  Its unnormalized update is
\[
    q_t(x_t,x_{t-1};\beta)
    =
    \widehat p_{t-1}(x_{t-1};\beta)
    f_\beta(x_t\mid x_{t-1})
    g_\beta(y_t\mid x_t).
\]
If one also wants to carry a learned static parameter as a random coordinate,
that coordinate can be appended to the state; the derivative below is still
with respect to the external \(\beta\), not with respect to an integration
variable.

Let \(z\in[-1,1]^D\) be reference coordinates and
\[
    r=\Psi_t(z)=(x_t,x_{t-1}).
\]
Let \(J_t(z)=|\det\nabla\Psi_t(z)|\).  The transformed target is
\[
    \widetilde q_t(z;\beta)
    =
    \widehat p_{t-1}(x_{t-1};\beta)
    f_\beta(x_t\mid x_{t-1})
    g_\beta(y_t\mid x_t)
    J_t(z).
    \tag{FB1}
\]
The branch fixes points \(z^{(j)}\), weights \(w_j\), basis functions, ranks,
and a core-construction algorithm.  The fitted square-root TT is
\[
    \phi_t(z;\beta)=G_{t,1}(z_1;\beta)\cdots G_{t,D}(z_D;\beta).
    \tag{FB2}
\]
The fixed-branch approximate density is
\[
    \widehat q_t(z;\beta)
    =
    e^{-c_t}\phi_t(z;\beta)^2+\tau_t\lambda_t(z),
    \qquad
    \int\lambda_t(z)\dd z=1.
    \tag{FB3}
\]
Its normalizer is
\[
    \widehat Z_t(\beta)
    =
    \int\widehat q_t(z;\beta)\dd z
    =
    e^{-c_t}\int\phi_t(z;\beta)^2\dd z+\tau_t.
    \tag{FB4}
\]
This is exactly the boxed convention (C11).
The next filter object is the normalized marginal
\[
    \widehat p_t(x_t;\beta)
    =
    \frac{
    \int \widehat q_t(z_t,z_{t-1};\beta)\dd z_{t-1}}
    {\widehat Z_t(\beta)}
    \left|\det\nabla z_t/\nabla x_t\right|.
    \tag{FB5}
\]
The Jacobian factor appears if the filter is evaluated in physical coordinates.
If the whole recursion is implemented in reference coordinates, it is absorbed
into the stored domain maps.

\begin{proposition}[Fixed-branch recursion is a normalized approximate filter]
Fix a time horizon \(T\), model densities \(f_\beta,g_\beta\), initial density
\(p_\beta(x_0)\), and fixed branch objects for each time step.  Suppose every
\(\widehat q_t(\cdot;\beta)\) defined by (FB3) is nonnegative, integrable, and
has \(0<\widehat Z_t(\beta)<\infty\).  Define \(\widehat p_t\) by (FB5),
starting from the fixed initial approximation.  Then each \(\widehat p_t\) is a
normalized density on the retained state variable \(x_t\).  Moreover, the
recursion is exactly the Bayes filtering recursion applied to the declared
approximate joint density \(\widehat q_t\), not to the exact joint density.
\end{proposition}

\begin{proof}
Nonnegativity follows from
\(\phi_t^2\ge0\), \(\tau_t>0\), and \(\lambda_t\ge0\).  Since
\(\widehat Z_t=\int\widehat q_t\) is positive and finite, the normalized joint
density
\[
    \widehat p_t^{\rm joint}(z)
    =
    \widehat q_t(z)/\widehat Z_t
\]
integrates to one.  The retained filter is the marginal of this normalized
joint density:
\[
    \int \widehat p_t(z_t)\dd z_t
    =
    \int
    \left[
    \int \widehat p_t^{\rm joint}(z_t,z_{t-1})\dd z_{t-1}
    \right]\dd z_t
    =
    \int \widehat p_t^{\rm joint}(z)\dd z
    =
    1.
\]
The change-of-variables factor in physical coordinates is the standard
Jacobian factor for density transformation, so normalization is preserved.  The
construction uses the exact Bayes filtering algebra with \(q_t\) replaced by
the declared approximation \(\widehat q_t\).  No step asserts that
\(\widehat q_t=q_t\), so the result is a normalized approximate filtering
recursion rather than an exact-posterior theorem.
\end{proof}


\section{Integrated Fixed-Branch Gradient Expansion}

The preceding sections preserved the source-order Zhao--Cui annotation and the
BayesFilter fixed-branch setup.  The next part expands the derivative path in
small mathematical steps.  It deliberately starts again from scalar
normalizers before returning to tensor trains, because the key claim is that
the gradient differentiates the same approximate scalar computed by the fixed
branch.

The merge has one important reading convention.  The Zhao--Cui annotation uses
\(\theta\) for a random static parameter inside the sequential posterior when
the paper studies joint state and parameter learning.  The fixed-branch
gradient part uses \(\beta\) for an external parameter value at which
BayesFilter evaluates an approximate likelihood.  These are different roles:
\[
    \theta \text{ can be an integration coordinate in Zhao--Cui,}
    \qquad
    \beta \text{ is the differentiated argument of } \widehat\ell_T .
    \tag{P20-B1}
\]
The fixed-branch scalar therefore does not claim to differentiate the adaptive
Zhao--Cui algorithm with respect to every internal coordinate.  It freezes the
realized numerical branch and differentiates the scalar produced by that
branch as \(\beta\) varies:
\[
    B=\{\text{domains, points, basis, ranks, shifts, sweeps, tolerances}\},
    \qquad
    \beta\mapsto \widehat\ell_T(\beta;B).
    \tag{P20-B2}
\]
This distinction is what makes the next derivation useful rather than merely
formal.  Without it, changing \(\beta\) could also change ranks, pivots, maps,
or fitting points; then two nearby likelihood evaluations would not be values
of one mathematical function.  With it, the fitted core values still change
with \(\beta\), but they change through the same stored fitting equations:
\[
    C_{t,k}(\beta;B)
    =
    \text{output of the fixed core-fitting rule at parameter }\beta.
    \tag{P20-B3}
\]
Thus the derivative below is neither a derivative of frozen numbers nor a
derivative of an adaptive search procedure.  It is the derivative of the
declared fixed-branch approximation.  This is the contract a finite-difference
diagnostic and a later implementation must obey.

\section{Fixed-Branch Gradient Motivation}

The filtering algorithm produces, at each time \(t\), an approximate
normalizing constant
\[
    \widehat Z_t(\beta)>0.
    \tag{P19-1}
\]
The parameter \(\beta\) is the model parameter at which we evaluate the
approximate likelihood.  The total approximate log likelihood is
\[
    \widehatell_T(\beta)
    =
    \sum_{t=1}^{T}\log \widehat Z_t(\beta).
    \tag{P19-2}
\]
Formally, \(\widehat Z_t(\beta)\) plays the role of a
partition-function-like normalizer: it is an integral of a nonnegative weight.
Here
\(\widehat Z_t(\beta)\) is also an integral of a nonnegative weight:
\[
    \widehat Z_t(\beta)
    =
    \int \widehat q_t(z;\beta)\dd z.
    \tag{P19-3}
\]
The variable \(z\) is a numerical coordinate, usually a transformed coordinate
on a fixed box.  The function \(\widehat q_t\) is not the exact Bayesian
density.  It is the squared tensor-train approximation used by the numerical
filter.

The panel cares about the gradient because parameter inference algorithms need
to know how the approximate log likelihood changes when \(\beta\) changes:
\[
    \nabla_\beta \widehatell_T(\beta).
    \tag{P19-4}
\]
The delicate point is not the derivative of \(\log\).  The delicate point is
that \(\widehat Z_t\) is produced by a complicated approximation algorithm.  If
the approximation algorithm changes its grid, ranks, pivots, or scaling shifts
when \(\beta\) changes, then the two likelihood evaluations may not be
evaluations of the same mathematical formula.  A classical analytical
derivative can only differentiate one declared formula.

That is why this note uses a fixed branch.

\begin{center}
\fbox{\begin{minipage}{0.92\linewidth}
\textbf{Fixed branch.}  Before differentiating, freeze the computational
choices: coordinate domain, fitting points, basis functions, tensor-train
ranks, sweep order, ridge parameter, defensive density, stabilizing shifts, and
the rule for constructing every fitted core.  Do not freeze the fitted core
values themselves.  Those core values are outputs of the fixed fitting
equations, so they remain functions of \(\beta\).  Then differentiate the
scalar defined by the saved choices and by the core values obtained from the
same saved fitting equations.
\end{minipage}}
\end{center}

The goal is therefore:
\[
    \text{differentiate the scalar actually computed by the saved forward pass.}
    \tag{P19-5}
\]
The goal is not:
\[
    \text{differentiate every possible adaptive version of the algorithm.}
    \tag{P19-6}
\]

\section{Warmup 1: A Normalizing Constant}

Start with the simplest possible object.  Let
\[
    Z(\beta)=\int_\Omega q(z;\beta)\dd z,
    \qquad
    q(z;\beta)\ge 0,
    \qquad
    0<Z(\beta)<\infty.
    \tag{P19-7}
\]
The normalized density is
\[
    p(z;\beta)=\frac{q(z;\beta)}{Z(\beta)}.
    \tag{P19-8}
\]
The log normalizing constant is
\[
    L(\beta)=\log Z(\beta).
    \tag{P19-9}
\]
Assume that \(q(z;\beta)\) is differentiable in \(\beta\), and that we may move
the derivative inside the integral:
\[
    \frac{\partial}{\partial\beta}
    \int_\Omega q(z;\beta)\dd z
    =
    \int_\Omega
    \frac{\partial q(z;\beta)}{\partial\beta}\dd z.
    \tag{P19-10}
\]
A sufficient condition is that, near the parameter value of interest, the
absolute derivative is bounded by an integrable function:
\[
    \left|
    \frac{\partial q(z;\beta)}{\partial\beta}
    \right|
    \le m(z),
    \qquad
    \int_\Omega m(z)\dd z<\infty.
    \tag{P19-11}
\]
Then the derivative of the log normalizer is
\[
\begin{aligned}
    \frac{\partial L(\beta)}{\partial\beta}
    &=
    \frac{\partial}{\partial\beta}\log Z(\beta)\\
    &=
    \frac{1}{Z(\beta)}
    \frac{\partial Z(\beta)}{\partial\beta}\\
    &=
    \frac{1}{Z(\beta)}
    \int_\Omega
    \frac{\partial q(z;\beta)}{\partial\beta}\dd z.
\end{aligned}
\tag{P19-12}
\]
This is the first idea in the whole gradient.  Everything later is only a way
to compute the numerator and denominator in (12) for the tensor-train
approximation.

For the sequential filter, the same formula is used at every time:
\[
    \frac{\partial\widehatell_T(\beta)}{\partial\beta}
    =
    \sum_{t=1}^{T}
    \frac{1}{\widehat Z_t(\beta)}
    \frac{\partial\widehat Z_t(\beta)}{\partial\beta}.
    \tag{P19-13}
\]

\section{Warmup 2: A Squared Approximation}

The fixed-branch squared tensor-train approximation has the form
\[
    \widehat q(z;\beta)
    =
    e^{-c}\phi(z;\beta)^2+\tau\lambda(z).
    \tag{P19-14}
\]
Here:
\[
\begin{array}{ll}
c & \text{is a frozen stabilizing shift},\\
\tau>0 & \text{is a frozen defensive mass},\\
\lambda(z)\ge0 & \text{is a frozen reference density with }\int\lambda=1,\\
\phi(z;\beta) & \text{is the fitted square-root function.}
\end{array}
\tag{P19-15}
\]
Because \(c,\tau,\lambda\) are frozen, the only object in (14) that changes
with \(\beta\) is \(\phi\).  Differentiating (14) gives
\[
\begin{aligned}
    \dot{\widehat q}(z;\beta)
    &=
    \frac{\partial}{\partial\beta}
    \left[
    e^{-c}\phi(z;\beta)^2+\tau\lambda(z)
    \right]\\
    &=
    e^{-c}
    \frac{\partial}{\partial\beta}\phi(z;\beta)^2
    +
    \frac{\partial}{\partial\beta}\tau\lambda(z)\\
    &=
    2e^{-c}\phi(z;\beta)\dot\phi(z;\beta).
\end{aligned}
\tag{P19-16}
\]
The dot means \(\partial/\partial\beta\).

The corresponding normalizer is
\[
\begin{aligned}
    \widehat Z(\beta)
    &=
    \int \widehat q(z;\beta)\dd z\\
    &=
    e^{-c}\int \phi(z;\beta)^2\dd z
    +
    \tau\int\lambda(z)\dd z\\
    &=
    e^{-c}\int \phi(z;\beta)^2\dd z+\tau.
\end{aligned}
\tag{P19-17}
\]
Using (16), its derivative is
\[
\begin{aligned}
    \dot{\widehat Z}(\beta)
    &=
    \int \dot{\widehat q}(z;\beta)\dd z\\
    &=
    2e^{-c}
    \int \phi(z;\beta)\dot\phi(z;\beta)\dd z.
\end{aligned}
\tag{P19-18}
\]
Substitution into (12) gives
\[
    \frac{\partial}{\partial\beta}\log\widehat Z(\beta)
    =
    \frac{
    2e^{-c}\int \phi(z;\beta)\dot\phi(z;\beta)\dd z}
    {e^{-c}\int\phi(z;\beta)^2\dd z+\tau}.
    \tag{P19-19}
\]

If \(c\), \(\tau\), or \(\lambda\) changed with \(\beta\), then (16) would have
extra terms:
\[
    \dot{\widehat q}
    =
    -\dot c\,e^{-c}\phi^2
    +
    2e^{-c}\phi\dot\phi
    +
    \dot\tau\,\lambda
    +
    \tau\dot\lambda.
    \tag{P19-20}
\]
The fixed-branch derivative does not use (20).  It uses the fixed-branch
formula (16).  This is not a minor technicality.  It is the mathematical line
between differentiating one saved scalar and differentiating a changing
adaptive procedure.

\section{Warmup 3: Two Coordinates, Rank One}

Before a full tensor train, consider a two-coordinate rank-one function:
\[
    \phi(z_1,z_2;\beta)=h_1(z_1;\beta)h_2(z_2;\beta).
    \tag{P19-21}
\]
The square integral is
\[
\begin{aligned}
    R(\beta)
    &=
    \int\!\!\int
    \phi(z_1,z_2;\beta)^2\dd z_1\dd z_2\\
    &=
    \int\!\!\int
    h_1(z_1;\beta)^2h_2(z_2;\beta)^2\dd z_1\dd z_2\\
    &=
    \left[\int h_1(z_1;\beta)^2\dd z_1\right]
    \left[\int h_2(z_2;\beta)^2\dd z_2\right].
\end{aligned}
\tag{P19-22}
\]
Define
\[
    R_1(\beta)=\int h_1(z_1;\beta)^2\dd z_1,
    \qquad
    R_2(\beta)=\int h_2(z_2;\beta)^2\dd z_2.
    \tag{P19-23}
\]
Then
\[
    R(\beta)=R_1(\beta)R_2(\beta).
    \tag{P19-24}
\]
Differentiate by the product rule:
\[
    \dot R(\beta)=\dot R_1(\beta)R_2(\beta)+R_1(\beta)\dot R_2(\beta).
    \tag{P19-25}
\]
Each one-dimensional derivative is
\[
    \dot R_1(\beta)=2\int h_1(z_1;\beta)\dot h_1(z_1;\beta)\dd z_1,
    \tag{P19-26}
\]
\[
    \dot R_2(\beta)=2\int h_2(z_2;\beta)\dot h_2(z_2;\beta)\dd z_2.
    \tag{P19-27}
\]
Combining (25)--(27),
\[
\begin{aligned}
    \dot R(\beta)
    &=
    2
    \left[\int h_1\dot h_1\dd z_1\right]
    \left[\int h_2^2\dd z_2\right]\\
    &\quad+
    2
    \left[\int h_1^2\dd z_1\right]
    \left[\int h_2\dot h_2\dd z_2\right].
\end{aligned}
\tag{P19-28}
\]
This is the simplest version of the tensor-train mass derivative: when one
factor changes, keep the other factors fixed and sum the contributions.

\section{Warmup 4: Two Coordinates, Rank \texorpdfstring{\(R\)}{R}}

Now let
\[
    \phi(z_1,z_2;\beta)
    =
    \sum_{r=1}^{R}h_{1r}(z_1;\beta)h_{2r}(z_2;\beta).
    \tag{P19-29}
\]
The square integral is
\[
\begin{aligned}
    R(\beta)
    &=
    \int\!\!\int
    \left[
    \sum_{r=1}^{R}h_{1r}(z_1)h_{2r}(z_2)
    \right]
    \left[
    \sum_{s=1}^{R}h_{1s}(z_1)h_{2s}(z_2)
    \right]
    \dd z_1\dd z_2\\
    &=
    \sum_{r=1}^{R}\sum_{s=1}^{R}
    \left[\int h_{1r}(z_1)h_{1s}(z_1)\dd z_1\right]
    \left[\int h_{2r}(z_2)h_{2s}(z_2)\dd z_2\right].
\end{aligned}
\tag{P19-30}
\]
Define the two mass matrices
\[
    M_1[r,s]=\int h_{1r}(z_1)h_{1s}(z_1)\dd z_1,
    \tag{P19-31}
\]
\[
    M_2[r,s]=\int h_{2r}(z_2)h_{2s}(z_2)\dd z_2.
    \tag{P19-32}
\]
Then
\[
    R(\beta)=\sum_{r,s}M_1[r,s]M_2[r,s].
    \tag{P19-33}
\]
This is the Frobenius inner product
\[
    R(\beta)=\langle M_1(\beta),M_2(\beta)\rangle_F.
    \tag{P19-34}
\]
Differentiate:
\[
    \dot R
    =
    \langle \dot M_1,M_2\rangle_F
    +
    \langle M_1,\dot M_2\rangle_F.
    \tag{P19-35}
\]
The derivative of a mass matrix entry is
\[
    \dot M_1[r,s]
    =
    \int
    \dot h_{1r}(z_1)h_{1s}(z_1)
    +
    h_{1r}(z_1)\dot h_{1s}(z_1)
    \dd z_1,
    \tag{P19-36}
\]
and similarly for \(M_2\).  The full tensor-train mass recursion is only this
calculation repeated along a chain of coordinates.

\section{Warmup 5: A Fixed Linear Solve}

The fitted tensor-train cores are obtained from finite-dimensional least-squares
updates.  One such update can be written
\[
    N(\beta)g(\beta)=d(\beta).
    \tag{P19-37}
\]
Here \(g\) is the vector of unknown core coefficients in the current update,
\(N\) is the regularized normal matrix, and \(d\) is the right-hand side.
Differentiate both sides:
\[
    \frac{\partial}{\partial\beta}\{N(\beta)g(\beta)\}
    =
    \frac{\partial d(\beta)}{\partial\beta}.
    \tag{P19-38}
\]
Use the product rule:
\[
    \dot N(\beta)g(\beta)+N(\beta)\dot g(\beta)=\dot d(\beta).
    \tag{P19-39}
\]
Move the first term to the right:
\[
    N(\beta)\dot g(\beta)=\dot d(\beta)-\dot N(\beta)g(\beta).
    \tag{P19-40}
\]
If \(N(\beta)\) is nonsingular, then
\[
    \dot g(\beta)=
    N(\beta)^{-1}\left[\dot d(\beta)-\dot N(\beta)g(\beta)\right].
    \tag{P19-41}
\]
In computation, one should not form the inverse.  One solves the linear system
(40) using the same numerical linear algebra used for the forward solve.

The ridge term is important because the normal matrix has the form
\[
    N=A^\top W A+\rho I,
    \qquad \rho>0.
    \tag{P19-42}
\]
If \(W\) is positive semidefinite, then
\[
    v^\top N v
    =
    v^\top A^\top W A v+\rho v^\top v
    \ge
    \rho\|v\|^2
    \tag{P19-43}
\]
for all \(v\).  Hence \(N\) is positive definite and nonsingular.  Without
nonsingularity, the derivative of the solve may not be well defined.

\section{The Fixed-Branch Forward Pass}

At time \(t\), the fixed branch constructs a nonnegative approximation
\[
    \widehat q_t(z;\beta)
    =
    e^{-c_t}\phi_t(z;\beta)^2+\tau_t\lambda_t(z)
    \tag{P19-44}
\]
and its normalizer
\[
    \widehat Z_t(\beta)=\int \widehat q_t(z;\beta)\dd z.
    \tag{P19-45}
\]
The branch is the collection of structural choices that define the scalar as a
function of \(\beta\).  These choices include the domain, fitting points, basis,
ranks, sweep order, shift, defensive term, and ridge parameter.  The branch does
not mean that the fitted numerical core values are held constant when
\(\beta\) changes.  Instead,
\[
    C_{t,k}(\beta)
    =
    \text{the core obtained by the saved fitting rule at parameter }\beta.
    \tag{P19-45a}
\]
The derivative pass computes \(\dot C_{t,k}(\beta)\).  This distinction matters:
freezing the core values would usually make the normalizer nearly constant and
would not assess the derivative of the fitted approximation.

\begin{center}
\begin{tabular}{p{0.27\linewidth}p{0.31\linewidth}p{0.31\linewidth}}
\toprule
Forward object & What is stored & Why it is needed later\\
\midrule
Domain map \(\Psi_t\) & Coordinate intervals and Jacobian & Evaluates the same transformed target at every derivative point\\
Fitting points \(z^{(j)}\) & Deterministic point array and weights & Defines the least-squares equations\\
Basis \(b_k\) & Basis family, degree, mass matrices & Evaluates cores and mass contractions\\
Ranks \(R_k\) & Fixed rank vector & Fixes core dimensions\\
Sweep order & Core update order and number of sweeps & Defines a finite composition of solve maps\\
Shift \(c_t\) & One frozen scalar & Defines the same scaled square-root target\\
Defensive term & \(\tau_t,\lambda_t\) & Ensures nonnegative density floor\\
Cores \(C_{t,k}(\beta)\) & Parameter-dependent fitted coefficient arrays & Evaluate \(\phi_t\) and \(\widehat q_t\)\\
Mass contractions & \(M_{t,>k}(\beta)\), \(M_{t,<k}(\beta)\) & Compute integrals and marginals\\
Normalizer & \(\widehat Z_t\) & Adds \(\log\widehat Z_t\) to the scalar\\
Carried filter & Numerator \(a_t\) and normalized \(\widehat p_t\) & Feeds the next time step\\
\bottomrule
\end{tabular}
\end{center}

The transformed target fitted at time \(t\) is
\[
    \widetilde q_t(z;\beta)
    =
    \widehat p_{t-1}(x_{t-1}(z);\beta)
    f_\beta(x_t(z)\mid x_{t-1}(z))
    g_\beta(y_t\mid x_t(z))
    J_t(z).
    \tag{P19-46}
\]
The fitted square-root values are
\[
    y_j(\beta)
    =
    e^{c_t/2}\sqrt{\widetilde q_t(z^{(j)};\beta)}.
    \tag{P19-47}
\]
The tensor train is
\[
    \phi_t(z;\beta)
    =
    H_{t,1}(z_1;\beta)H_{t,2}(z_2;\beta)\cdots H_{t,D}(z_D;\beta),
    \tag{P19-48}
\]
where \(H_{t,k}(z_k)\in\R^{R_{k-1}\times R_k}\) and \(R_0=R_D=1\).

The retained numerator for the next filter is
\[
    a_t(z_t;\beta)
    =
    \int \widehat q_t(z_t,z_{t-1};\beta)\dd z_{t-1}.
    \tag{P19-49}
\]
The retained normalized filter is
\[
    \widehat p_t(z_t;\beta)
    =
    \frac{a_t(z_t;\beta)}{\widehat Z_t(\beta)}.
    \tag{P19-50}
\]
If the filter is reported in physical coordinates, multiply by the appropriate
change-of-variables factor.  The derivative logic is unchanged if the domain
map is frozen.

\section{The Fixed-Branch Derivative Pass}

The derivative pass computes the derivative of each forward object that depends
on \(\beta\).  It does not create a new branch.

\begin{center}
\begin{tabular}{p{0.32\linewidth}p{0.32\linewidth}p{0.26\linewidth}}
\toprule
Derivative object & Forward object differentiated & Formula source\\
\midrule
\(\dot{\widetilde q}_{t,j}\) & target value at fitting point & product rule\\
\(\dot y_j\) & square-root fit value & chain rule\\
\(\dot A_k,\dot N_k,\dot d_k\) & least-squares system & environment derivatives\\
\(\dot C_{t,k}\) & fitted core & fixed linear solve\\
\(\dot M_{t,>k}\) & mass contraction & product rule in core indices\\
\(\dot{\widehat Z}_t\) & normalizer & squared-integral derivative\\
\(\dot a_t\) & carried numerator & marginal contraction derivative\\
\(\dot{\widehat p}_t\) & normalized carried filter & quotient rule\\
\(\dot{\widetilde q}_{t+1}\) & next target & product rule using \(\dot{\widehat p}_t\)\\
\bottomrule
\end{tabular}
\end{center}

\subsection{Target And Square-Root Target}

Differentiate (46).  At fitting point \(z^{(j)}\), abbreviate
\[
    \widehat p_{t-1,j}=\widehat p_{t-1}(x_{t-1}^{(j)};\beta),
    \quad
    f_j=f_\beta(x_t^{(j)}\mid x_{t-1}^{(j)}),
    \quad
    g_j=g_\beta(y_t\mid x_t^{(j)}).
    \tag{P19-51}
\]
With \(J_{t,j}\) frozen, the derivative is
\[
\begin{aligned}
    \dot{\widetilde q}_{t,j}
    &=
    J_{t,j}
    \left[
    \dot{\widehat p}_{t-1,j} f_j g_j
    +
    \widehat p_{t-1,j}\dot f_j g_j
    +
    \widehat p_{t-1,j} f_j\dot g_j
    \right].
\end{aligned}
\tag{P19-52}
\]
If the domain map depended on \(\beta\), an additional
\(\widehat p f g\,\dot J\) term would appear.  That is outside the fixed branch
used here.

From (47),
\[
    \dot y_j
    =
    e^{c_t/2}
    \frac{1}{2\sqrt{\widetilde q_{t,j}}}
    \dot{\widetilde q}_{t,j},
    \tag{P19-53}
\]
provided \(\widetilde q_{t,j}>0\).  A numerical implementation should use a
positive floor if the target is close to zero.  If such a floor is used, the
flooring rule becomes part of the declared scalar.  For example,
\[
    y_j(\beta)
    =
    e^{c_t/2}\sqrt{\max\{\widetilde q_{t,j}(\beta),\varepsilon_{\rm floor}\}}
    \tag{P19-53a}
\]
has derivative zero through the floor whenever
\(\widetilde q_{t,j}<\varepsilon_{\rm floor}\).  The formulas below describe
the smooth positive-target case; a floored implementation must differentiate
the floored scalar it actually evaluates.

\subsection{Design Matrix And Core Solve}

The design row is the first place where tensor-train notation can look more
mysterious than it is.  Fix one sample point \(z^{(j)}\) and one core \(k\).
All cores except core \(k\) are treated as the known left and right
environments for this local least-squares update.  Write
\[
    L_{j,k}[a]
    =
    \left[
    H_{1}(z_1^{(j)})\cdots H_{k-1}(z_{k-1}^{(j)})
    \right]_{1a},
    \tag{P19-53b}
\]
\[
    R_{j,k}[b]
    =
    \left[
    H_{k+1}(z_{k+1}^{(j)})\cdots H_D(z_D^{(j)})
    \right]_{b1}.
    \tag{P19-53c}
\]
The local core evaluated at \(z_k^{(j)}\) is
\[
    H_k(z_k^{(j)})[a,b]
    =
    \sum_{\ell=1}^{m_k}
    C_k[a,\ell,b]\,b_k(z_k^{(j)})[\ell],
    \tag{P19-53d}
\]
where \(m_k\) is the number of one-dimensional basis functions in coordinate
\(k\).  The value of the tensor train at this one point is therefore
\[
\begin{aligned}
    \phi_t(z^{(j)};\beta)
    &=
    \sum_{a=1}^{R_{k-1}}
    \sum_{b=1}^{R_k}
    L_{j,k}[a]\,H_k(z_k^{(j)})[a,b]\,R_{j,k}[b]\\
    &=
    \sum_{a=1}^{R_{k-1}}
    \sum_{\ell=1}^{m_k}
    \sum_{b=1}^{R_k}
    L_{j,k}[a]\,
    b_k(z_k^{(j)})[\ell]\,
    R_{j,k}[b]\,
    C_k[a,\ell,b].
\end{aligned}
    \tag{P19-53e}
\]
Equation (53e) is the whole reason the local fit is a linear least-squares
problem.  At the fixed sample point, the unknown numbers are the entries of
the core \(C_k\).  Every coefficient multiplying \(C_k[a,\ell,b]\) is already
known from the branch and from the other cores.

If \(g_k=\vecop(C_k)\) stacks entries in the order \((a,\ell,b)\), the row
coefficient of \(g_k[a,\ell,b]\) is
\[
    A_{j,k}[a,\ell,b]
    =
    L_{j,k}[a]\,
    b_k(z_k^{(j)})[\ell]\,
    R_{j,k}[b].
    \tag{P19-53f}
\]
The compact Kronecker expression in (57) is just (53f) written without three
sums:
\[
    A_{j,k}
    =
    R_{j,k}^{\top}\otimes b_k(z_k^{(j)})^\top\otimes L_{j,k}.
    \tag{P19-53g}
\]
To differentiate the row, keep the basis value fixed and differentiate the
two environments:
\[
\begin{aligned}
    \dot A_{j,k}[a,\ell,b]
    &=
    \dot L_{j,k}[a]\,
    b_k(z_k^{(j)})[\ell]\,
    R_{j,k}[b]\\
    &\quad+
    L_{j,k}[a]\,
    b_k(z_k^{(j)})[\ell]\,
    \dot R_{j,k}[b].
\end{aligned}
    \tag{P19-53h}
\]
Equation (58) is (53h) written in the same Kronecker shorthand.  If the basis
or fitting point moved with \(\beta\), there would be an additional term with
\(\dot b_k(z_k^{(j)})\).  That term is absent because the branch keeps basis
and fitting points fixed.

For one core \(k\), the least-squares relation has rows
\[
    \phi_t(z^{(j)};\beta)=A_{j,k}(\beta)g_k(\beta),
    \tag{P19-54}
\]
where \(g_k=\vecop(C_{t,k})\).  Let
\[
    L_{j,k}
    =
    H_{1}(z_1^{(j)})\cdots H_{k-1}(z_{k-1}^{(j)}),
    \tag{P19-55}
\]
\[
    R_{j,k}
    =
    H_{k+1}(z_{k+1}^{(j)})\cdots H_D(z_D^{(j)}).
    \tag{P19-56}
\]
Then
\[
    A_{j,k}
    =
    R_{j,k}^{\top}\otimes b_k(z_k^{(j)})^\top\otimes L_{j,k}.
    \tag{P19-57}
\]
Because basis values and fitting points are frozen,
\[
    \dot A_{j,k}
    =
    \dot R_{j,k}^{\top}\otimes b_k(z_k^{(j)})^\top\otimes L_{j,k}
    +
    R_{j,k}^{\top}\otimes b_k(z_k^{(j)})^\top\otimes \dot L_{j,k}.
    \tag{P19-58}
\]
The environments are differentiated by product rule:
\[
    \dot L_{j,1}=0,
    \qquad
    \dot L_{j,k+1}
    =
    \dot L_{j,k}H_k(z_k^{(j)})
    +
    L_{j,k}\dot H_k(z_k^{(j)}),
    \tag{P19-59}
\]
\[
    \dot R_{j,D}=0,
    \qquad
    \dot R_{j,k-1}
    =
    \dot H_k(z_k^{(j)})R_{j,k}
    +
    H_k(z_k^{(j)})\dot R_{j,k}.
    \tag{P19-60}
\]
The local core matrix derivative is
\[
    \dot H_k(z_k^{(j)})[a,b]
    =
    \sum_{\ell}
    \dot C_k[a,\ell,b]\,b_k(z_k^{(j)})[\ell].
    \tag{P19-61}
\]

The fixed ridge solve is
\[
    N_k g_k=d_k,
    \qquad
    N_k=A_k^\top W A_k+\rho I,
    \qquad
    d_k=A_k^\top W y.
    \tag{P19-62}
\]
Using (40),
\[
    N_k\dot g_k=\dot d_k-\dot N_k g_k,
    \tag{P19-63}
\]
where
\[
    \dot N_k=\dot A_k^\top W A_k+A_k^\top W\dot A_k,
    \tag{P19-64}
\]
\[
    \dot d_k=\dot A_k^\top W y+A_k^\top W\dot y.
    \tag{P19-65}
\]
The derivative update is therefore another linear solve with the same
left-hand matrix \(N_k\).  This is why the fixed branch is implementable: the
forward pass and derivative pass use the same stored computational path.

\subsection{Mass Contraction And Normalizer}

The full right-mass recursion is the same idea as the rank-\(R\) warmup in
(29)--(36).  In the warmup, \(M_1[r,s]\) stored the inner products of the
first-coordinate factors and \(M_2[r,s]\) stored the inner products of the
second-coordinate factors.  In the full tensor train, \(M_{>j}[b,b']\) plays
the role of the already-integrated right factor, \(B_j[\ell,\ell']\) is the
one-coordinate basis inner product, and the two copies of \(C_j\) are the two
factors coming from squaring \(\phi_t\).

The right mass contraction is
\[
    M_{>j-1}[a,a']
    =
    \sum_{b,b',\ell,\ell'}
    C_j[a,\ell,b]
    M_{>j}[b,b']
    C_j[a',\ell',b']
    B_j[\ell,\ell'].
    \tag{P19-66}
\]
Here \(B_j\) is the basis mass matrix.  It is frozen.  Differentiate (66):
\[
\begin{aligned}
    \dot M_{>j-1}[a,a']
    &=
    \sum_{b,b',\ell,\ell'}
    \dot C_j[a,\ell,b]
    M_{>j}[b,b']
    C_j[a',\ell',b']
    B_j[\ell,\ell']\\
    &\quad+
    \sum_{b,b',\ell,\ell'}
    C_j[a,\ell,b]
    \dot M_{>j}[b,b']
    C_j[a',\ell',b']
    B_j[\ell,\ell']\\
    &\quad+
    \sum_{b,b',\ell,\ell'}
    C_j[a,\ell,b]
    M_{>j}[b,b']
    \dot C_j[a',\ell',b']
    B_j[\ell,\ell'].
\end{aligned}
\tag{P19-67}
\]
Start from
\[
    M_{>D}=1,
    \qquad
    \dot M_{>D}=0.
    \tag{P19-68}
\]
Then (66)--(67) move from the right end of the chain to the left end.  The
final square integral is
\[
    R_t(\beta)=\int\phi_t(z;\beta)^2\dd z=M_{>0}.
    \tag{P19-69}
\]
Its derivative is
\[
    \dot R_t(\beta)=\dot M_{>0}.
    \tag{P19-70}
\]
By (17)--(18),
\[
    \widehat Z_t(\beta)=e^{-c_t}R_t(\beta)+\tau_t,
    \tag{P19-71}
\]
\[
    \dot{\widehat Z}_t(\beta)=e^{-c_t}\dot R_t(\beta).
    \tag{P19-72}
\]

\subsection{Carried Filter And Next Time Step}

The carried numerator is
\[
    a_t(z_t;\beta)
    =
    \int \widehat q_t(z_t,z_{t-1};\beta)\dd z_{t-1}.
    \tag{P19-73}
\]
Its derivative is
\[
    \dot a_t(z_t;\beta)
    =
    \int \dot{\widehat q}_t(z_t,z_{t-1};\beta)\dd z_{t-1}.
    \tag{P19-74}
\]
For the squared part of the approximation, this marginal is another tensor
contraction.  To make the formula concrete, suppose the retained coordinate is
the last coordinate \(z_D=z_t\), and the previous-state coordinates are
\(z_1,\ldots,z_{D-1}\).  First integrate out the previous-state coordinates by
the same left-to-right mass recursion:
\[
    M_{<0}=1,
    \qquad
    M_{<j}[b,b']
    =
    \sum_{a,a',\ell,\ell'}
    M_{<j-1}[a,a']\,
    C_j[a,\ell,b]\,
    C_j[a',\ell',b']\,
    B_j[\ell,\ell']
    \tag{P19-74a}
\]
for \(j=1,\ldots,D-1\).  The derivative of this marginal environment is
\[
\begin{aligned}
    \dot M_{<j}[b,b']
    &=
    \sum_{a,a',\ell,\ell'}
    \dot M_{<j-1}[a,a']\,
    C_j[a,\ell,b]\,
    C_j[a',\ell',b']\,
    B_j[\ell,\ell']\\
    &\quad+
    \sum_{a,a',\ell,\ell'}
    M_{<j-1}[a,a']\,
    \dot C_j[a,\ell,b]\,
    C_j[a',\ell',b']\,
    B_j[\ell,\ell']\\
    &\quad+
    \sum_{a,a',\ell,\ell'}
    M_{<j-1}[a,a']\,
    C_j[a,\ell,b]\,
    \dot C_j[a',\ell',b']\,
    B_j[\ell,\ell'].
\end{aligned}
    \tag{P19-74b}
\]
Start with \(\dot M_{<0}=0\).  After the recursion reaches \(D-1\), the
retained-coordinate numerator is the function
\[
\begin{aligned}
    a_t(z_D;\beta)
    &=
    e^{-c_t}
    \sum_{b,b',\ell,\ell'}
    M_{<D-1}[b,b']\,
    C_D[b,\ell,1]\,
    C_D[b',\ell',1]\,
    b_D(z_D)[\ell]\,
    b_D(z_D)[\ell']\\
    &\quad+
    \tau_t
    \int \lambda_t(z_D,z_{1:D-1})\dd z_{1:D-1}.
\end{aligned}
    \tag{P19-74c}
\]
Its derivative, with \(c_t,\tau_t,\lambda_t\) and the basis fixed, is
\[
\begin{aligned}
    \dot a_t(z_D;\beta)
    &=
    e^{-c_t}
    \sum_{b,b',\ell,\ell'}
    \dot M_{<D-1}[b,b']\,
    C_D[b,\ell,1]\,
    C_D[b',\ell',1]\,
    b_D(z_D)[\ell]\,
    b_D(z_D)[\ell']\\
    &\quad+
    e^{-c_t}
    \sum_{b,b',\ell,\ell'}
    M_{<D-1}[b,b']\,
    \dot C_D[b,\ell,1]\,
    C_D[b',\ell',1]\,
    b_D(z_D)[\ell]\,
    b_D(z_D)[\ell']\\
    &\quad+
    e^{-c_t}
    \sum_{b,b',\ell,\ell'}
    M_{<D-1}[b,b']\,
    C_D[b,\ell,1]\,
    \dot C_D[b',\ell',1]\,
    b_D(z_D)[\ell]\,
    b_D(z_D)[\ell'].
\end{aligned}
    \tag{P19-74d}
\]
Equations (74a)--(74d) are the coding recipe behind the short integral
identity (74).  If the retained state is not the last coordinate, reorder the
coordinates or use left and right environments around the retained block; the
same product-rule pattern applies.

The normalized carried filter is
\[
    \widehat p_t(z_t;\beta)
    =
    \frac{a_t(z_t;\beta)}{\widehat Z_t(\beta)}.
    \tag{P19-75}
\]
Differentiate by quotient rule:
\[
    \dot{\widehat p}_t(z_t;\beta)
    =
    \frac{\dot a_t(z_t;\beta)}{\widehat Z_t(\beta)}
    -
    \frac{a_t(z_t;\beta)\dot{\widehat Z}_t(\beta)}
    {\widehat Z_t(\beta)^2}.
    \tag{P19-76}
\]
This is the crucial bridge between time steps.  At time \(t+1\), the target
contains \(\widehat p_t\), so its derivative contains
\(\dot{\widehat p}_t\).  Without saving the carried-filter derivative, the
next-step gradient is incomplete.

\subsection{Concrete One-Coordinate Carried-Filter Storage}

The quotient formula (76) is a mathematical identity.  A later implementation
also needs a concrete stored object.  In the minimal two-coordinate fixed
branch, use the normalized Legendre basis \(b_1(z)\in\R^p\) on the retained
coordinate and store the carried numerator as a coefficient matrix:
\[
    Q_t[\ell,m]
    =
    e^{-c_t}\sum_{r,s=1}^{R}
    C_1[0,\ell,r]S_2[r,s]C_1[0,m,s]
    +\tau_t\ind_{\ell=0}\ind_{m=0},
    \qquad
    \shape(Q_t)=(p,p).
    \tag{P22-K1}
\]
The retained numerator is then the quadratic form
\[
    a_t(z_1;\beta)=b_1(z_1)^\top Q_t b_1(z_1).
    \tag{P22-K2}
\]
The defensive term in (K1) gives \(\tau_t/2\) because the normalized Legendre
constant satisfies \(b_1(z)[0]=1/\sqrt2\), hence
\[
    \tau_t b_1(z_1)[0]^2=\tau_t/2.
    \tag{P22-K3}
\]
The derivative coefficient matrix is
\[
\begin{aligned}
    \dot Q_t[\ell,m]
    &=
    e^{-c_t}\sum_{r,s}
    \dot C_1[0,\ell,r]S_2[r,s]C_1[0,m,s]\\
    &\quad+
    e^{-c_t}\sum_{r,s}
    C_1[0,\ell,r]\dot S_2[r,s]C_1[0,m,s]\\
    &\quad+
    e^{-c_t}\sum_{r,s}
    C_1[0,\ell,r]S_2[r,s]\dot C_1[0,m,s],
    \qquad
    \shape(\dot Q_t)=(p,p).
\end{aligned}
\tag{P22-K4}
\]
The normalized carried filter is stored as
\[
    P_t=\frac{Q_t}{\widehat Z_t},
    \qquad
    \dot P_t=
    \frac{\dot Q_t}{\widehat Z_t}
    -
    \frac{Q_t\dot{\widehat Z}_t}{\widehat Z_t^2},
    \qquad
    \shape(P_t)=\shape(\dot P_t)=(p,p).
    \tag{P22-K5}
\]
For a query vector \(z^{\rm query}\in[-1,1]^M\), build
\[
    B^{\rm query}[j,\ell]=b_1(z^{\rm query}_j)[\ell],
    \qquad
    \shape(B^{\rm query})=(M,p).
    \tag{P22-K6}
\]
The carried evaluator returns
\[
\begin{aligned}
    \widehat p_t^{\rm ref}[j]
    &=
    \sum_{\ell,m}
    B^{\rm query}[j,\ell]P_t[\ell,m]B^{\rm query}[j,m],\\
    \dot{\widehat p}_t^{\rm ref}[j]
    &=
    \sum_{\ell,m}
    B^{\rm query}[j,\ell]\dot P_t[\ell,m]B^{\rm query}[j,m],
\end{aligned}
\qquad
\shape(\widehat p_t^{\rm ref})
=
\shape(\dot{\widehat p}_t^{\rm ref})
=(M,).
\tag{P22-K7}
\]
At the next time step, the query points are the previous-state coordinate of
the fitting table:
\[
    z^{\rm query}_j=Z_{\rm fit}[j,2],
    \qquad
    p_j=\widehat p_{t-1}^{\rm ref}[j],
    \qquad
    \dot p_j=\dot{\widehat p}_{t-1}^{\rm ref}[j].
    \tag{P22-K8}
\]
The stored object is therefore not an informal function name.  It is a pair of
coefficient matrices and a declared query rule.  This is what makes the
next-step target derivative mechanically reproducible.

\section{Proposition 1: The Fixed-Branch Recursion Is A Normalized Approximate Filter}

\begin{proposition}
Fix a parameter value \(\beta\), a time horizon \(T\), and fixed branch objects
for all times.  Suppose that each approximate joint density
\[
    \widehat q_t(z_t,z_{t-1};\beta)
    =
    e^{-c_t}\phi_t(z_t,z_{t-1};\beta)^2
    +
    \tau_t\lambda_t(z_t,z_{t-1})
    \tag{P19-77}
\]
is nonnegative and integrable, and that
\[
    0<\widehat Z_t(\beta)
    =
    \int \widehat q_t(z_t,z_{t-1};\beta)\dd z_t\dd z_{t-1}
    <\infty.
    \tag{P19-78}
\]
Define
\[
    \widehat p_t(z_t;\beta)
    =
    \frac{\int \widehat q_t(z_t,z_{t-1};\beta)\dd z_{t-1}}
    {\widehat Z_t(\beta)}.
    \tag{P19-79}
\]
Then \(\widehat p_t(\cdot;\beta)\) is a normalized density in \(z_t\).  The
recursion is a Bayesian filtering recursion for the declared approximate joint
density \(\widehat q_t\), not a proof that \(\widehat q_t\) equals the exact
Bayesian joint density.
\end{proposition}

\begin{proof}
Nonnegativity follows immediately:
\[
    e^{-c_t}\phi_t^2\ge0,
    \qquad
    \tau_t\lambda_t\ge0.
    \tag{P19-80}
\]
Therefore \(\widehat q_t\ge0\).  By assumption,
\(\widehat Z_t\) is positive and finite, so
\[
    \widehat p_t^{\rm joint}(z_t,z_{t-1};\beta)
    =
    \frac{\widehat q_t(z_t,z_{t-1};\beta)}
    {\widehat Z_t(\beta)}
    \tag{P19-81}
\]
is a normalized joint density:
\[
\begin{aligned}
    \int\!\!\int
    \widehat p_t^{\rm joint}(z_t,z_{t-1};\beta)
    \dd z_t\dd z_{t-1}
    &=
    \frac{1}{\widehat Z_t(\beta)}
    \int\!\!\int
    \widehat q_t(z_t,z_{t-1};\beta)
    \dd z_t\dd z_{t-1}\\
    &=
    1.
\end{aligned}
\tag{P19-82}
\]
The retained filter is the marginal of this normalized joint density:
\[
    \widehat p_t(z_t;\beta)
    =
    \int
    \widehat p_t^{\rm joint}(z_t,z_{t-1};\beta)
    \dd z_{t-1}.
    \tag{P19-83}
\]
Integrating (83) over \(z_t\),
\[
\begin{aligned}
    \int \widehat p_t(z_t;\beta)\dd z_t
    &=
    \int
    \left[
    \int \widehat p_t^{\rm joint}(z_t,z_{t-1};\beta)
    \dd z_{t-1}
    \right]\dd z_t\\
    &=
    \int\!\!\int
    \widehat p_t^{\rm joint}(z_t,z_{t-1};\beta)
    \dd z_t\dd z_{t-1}\\
    &=
    1.
\end{aligned}
\tag{P19-84}
\]
Thus \(\widehat p_t\) is normalized.  The proof used only the declared
approximation \(\widehat q_t\).  It did not use the exact Bayesian density
except as motivation for constructing \(\widehat q_t\).  Hence the theorem is a
normalization theorem for the approximate filtering recursion.
\end{proof}

\section{Proposition 2: The Fixed-Branch Derivative Differentiates The Same Scalar}

\begin{lemma}[Fixed ridge sweeps are differentiable maps]
Fix one branch: the coordinate domains, basis functions, fitting points,
weights, ranks, sweep order, ridge parameter, defensive terms, and shifts are
all held fixed.  Consider one core update written as a ridge least-squares
normal equation
\[
    N_k(\beta)g_k(\beta)=d_k(\beta),
    \qquad
    N_k(\beta)=A_k(\beta)^\top W_k A_k(\beta)+\varepsilon_k I,
\]
where \(\varepsilon_k>0\), \(W_k\) is fixed positive semidefinite, and
\(N_k(\beta)\) is nonsingular in the neighborhood being differentiated.  If the
entries of \(A_k(\beta)\) and \(d_k(\beta)\) are differentiable, then
\[
    \dot g_k(\beta)
    =
    N_k(\beta)^{-1}
    \left[
    \dot d_k(\beta)-\dot N_k(\beta)g_k(\beta)
    \right].
\]
Thus this fixed core update is differentiable.  A fixed finite sweep is a
finite composition of such updates, so the fitted cores produced by the fixed
sweep rule are differentiable functions of \(\beta\).
\end{lemma}

\begin{proof}
Because \(N_k(\beta)\) is nonsingular, the core vector is
\[
    g_k(\beta)=N_k(\beta)^{-1}d_k(\beta).
\]
Equivalently, it satisfies \(N_k g_k=d_k\).  Differentiate this identity:
\[
    \dot N_k g_k+N_k\dot g_k=\dot d_k.
\]
Solving for \(\dot g_k\) gives the displayed formula.  The assumptions that the
branch, points, weights, ranks, ridge parameter, shifts, and sweep order are
fixed ensure that no discrete choice changes while differentiating.  Therefore
each update is an ordinary differentiable map, and composing finitely many
updates gives an ordinary differentiable fitted-core map.
\end{proof}

\begin{proposition}
Fix the branch \(B=(B_1,\ldots,B_T)\).  Assume:
\[
    \widehatell_T(\beta;B)=\sum_{t=1}^T\log\widehat Z_t(\beta;B_t),
    \tag{P19-85}
\]
each \(\widehat Z_t\) is positive and finite, all model density evaluations are
differentiable in \(\beta\), every least-squares update uses a fixed
nonsingular ridge system, and differentiation under the finite-dimensional
contractions and integrals is valid.  Then the derivative pass defined by
(52)--(76) computes
\[
    \frac{\partial}{\partial\beta}
    \widehatell_T(\beta;B).
    \tag{P19-86}
\]
\end{proposition}

\begin{proof}
The proof follows the computational graph of the fixed branch.

First, the transformed target (46) is a product of the previous carried filter,
transition density, observation density, and frozen Jacobian.  The product rule
therefore gives (52).  If \(t=1\), the previous filter is the initial density;
for \(t>1\), its derivative is supplied by the previous carried-filter
derivative (76).

Second, the fitted target values (47) are differentiable functions of
\(\widetilde q_{t,j}\).  The chain rule gives (53).  The stabilizing shift
\(c_t\) is frozen, so no \(\dot c_t\) term appears.

Third, each core update is a fixed linear solve.  Equations (37)--(41) prove
that differentiating the solve yields
\[
    N_k\dot g_k=\dot d_k-\dot N_k g_k.
    \tag{P19-87}
\]
Equations (54)--(65) identify \(N_k,d_k,\dot N_k,\dot d_k\) for the fixed
least-squares update.  Since the number and order of solves are fixed, a finite
composition of differentiable solve maps gives differentiable fitted cores.

Fourth, the square integral \(R_t=\int\phi_t^2\) is computed by the mass
recursion (66).  Its derivative is the product-rule recursion (67).  Therefore
\[
    \dot R_t=\frac{\partial}{\partial\beta}
    \int\phi_t(z;\beta)^2\dd z.
    \tag{P19-88}
\]
With frozen \(c_t,\tau_t,\lambda_t\), (71)--(72) give the derivative of
\(\widehat Z_t\).

Fifth, the carried numerator and carried filter are defined by (73)--(75).
Their derivatives are (74) and (76), obtained by differentiating under the
integral and by the quotient rule.  This supplies the derivative object needed
by the next time step.

Finally, since each time contribution is \(\log\widehat Z_t(\beta;B_t)\),
Warmup 1 gives
\[
    \frac{\partial}{\partial\beta}
    \log\widehat Z_t(\beta;B_t)
    =
    \frac{\dot{\widehat Z}_t(\beta;B_t)}
    {\widehat Z_t(\beta;B_t)}.
    \tag{P19-89}
\]
Summing (89) over \(t=1,\ldots,T\) gives (86).

Every object differentiated in this proof is an object in the saved branch
\(B\).  The proof never differentiates a changed rank, changed grid, changed
shift, changed pivot set, or changed domain.  Therefore the derivative is the
derivative of the same scalar \(\widehatell_T(\beta;B)\), not the derivative of
an adaptive algorithm that chooses a new \(B\) when \(\beta\) changes.
\end{proof}

\section{What This Does Not Prove}

The result above is deliberately narrow.

It does not prove exact posterior accuracy:
\[
    \widehat q_t=q_t
    \quad\text{is not claimed.}
    \tag{P19-90}
\]
It proves normalization and differentiation of the declared approximation.

It does not prove ranks stay small:
\[
    R_k\ \text{moderate for all }k
    \quad\text{is a diagnostic requirement, not a theorem here.}
    \tag{P19-91}
\]
It does not differentiate adaptive decisions:
\[
    \frac{\partial}{\partial\beta}
    \{\text{chosen ranks, pivots, domains, shifts, fitting points}\}
    \quad\text{is not part of the formula.}
    \tag{P19-92}
\]
It does not prove HMC convergence:
\[
    \text{correct local gradient of an approximate scalar}
    \ne
    \text{posterior sampling guarantee.}
    \tag{P19-93}
\]

These caveats make the result more honest, not weaker.  The fixed-branch
gradient is a precise mathematical contract.  Empirical accuracy, rank
stability, and sampling performance must be assessed separately.

\section{Finite-Difference Diagnostic}

The finite-difference diagnostic checks that the value path and derivative path agree
for the same branch.  Choose a base parameter \(\beta_0\).  Build the branch
once and save the structural choices
\[
    \begin{aligned}
    B=\{&
    \text{domains, points, basis, ranks, sweeps, shifts,}\\
    &\text{defensive terms, ridge parameters,}\\
    &\text{deterministic initialization rules}\}.
    \end{aligned}
    \tag{P19-94}
\]
Then evaluate the same scalar at perturbed parameters while reusing those
structural choices and recomputing the parameter-dependent fitted values by the
same fixed equations:
\[
    \widehatell_T(\beta_0+h e_i;B),
    \qquad
    \widehatell_T(\beta_0-h e_i;B).
    \tag{P19-95}
\]
That means the domains, points, ranks, shifts, and sweep order are unchanged,
but the fitted core values
\[
    C_{t,k}(\beta_0+h e_i;B)
    \quad\text{and}\quad
    C_{t,k}(\beta_0-h e_i;B)
    \tag{P19-95a}
\]
are recomputed from the same saved fitting rule.  They are not copied from
\(\beta_0\).
The centered finite difference is
\[
    D_i(h)
    =
    \frac{
    \widehatell_T(\beta_0+h e_i;B)
    -
    \widehatell_T(\beta_0-h e_i;B)}
    {2h}.
    \tag{P19-96}
\]
For the one-parameter minimal case, the same field is
\[
    D(h)
    =
    \frac{
    \widehatell_2(\beta_0+h;B)
    -
    \widehatell_2(\beta_0-h;B)}
    {2h}.
    \tag{P22-FD0a}
\]
Compare it to the analytical derivative
\[
    G_i=\partial_i\widehatell_T(\beta_0;B).
    \tag{P19-97}
\]
The P22-local one-parameter derivative field is
\[
    G=\partial_\beta\widehatell_2(\beta_0;B).
    \tag{P22-FD0b}
\]
Record
\[
    e_i(h)=|D_i(h)-G_i|,
    \tag{P19-98}
\]
and
\[
    e_i^{\rm rel}(h)
    =
    \frac{|D_i(h)-G_i|}{1+|G_i|}.
    \tag{P19-99}
\]
For the one-parameter minimal case,
\[
    e(h)=|D(h)-G|,
    \qquad
    e_{\rm rel}(h)=\frac{|D(h)-G|}{1+|G|}.
    \tag{P22-FD0c}
\]
Use a small schedule such as
\[
    h\in\{10^{-2},10^{-3},10^{-4},10^{-5}\}.
    \tag{P19-100}
\]
The P22-local step-size field is
\[
    H_{\rm FD}=\{10^{-2},10^{-3},10^{-4},10^{-5}\}.
    \tag{P22-FD0d}
\]
The P22-local recompute-core rule is
\[
    C_{t,k}(\beta_0\pm h;B)
    =
    \text{the core produced by the fixed fitting rule at }\beta_0\pm h,
    \tag{P22-FD0e}
\]
not the core copied from \(\beta_0\).
The report should be a declared mathematical object, not an informal output
list.
For the one-parameter \(T=2\) minimal case, record
\[
\begin{aligned}
    \mathcal R_{\rm FD}
    =
    \{&
    \widehatell_2(\beta_0;B),\
    G,\
    (h_i,D(h_i),e(h_i),e_{\rm rel}(h_i))_{i=1}^{4},\\
    &I_-(h_i),I_+(h_i),K_-(h_i),K_+(h_i),\
    W_i,\ \text{status}\}.
\end{aligned}
\tag{P22-FD1}
\]
Here
\[
    I_\pm(h_i)
    =
    \ind\{B(\beta_0\pm h_i)=B\}
    \tag{P22-FD2}
\]
checks branch-manifest equality, and
\[
    K_\pm(h_i)
    =
    \ind\{C_k(\beta_0\pm h_i;B)
    \text{ was recomputed by the fixed fitting rule}\}
    \tag{P22-FD3}
\]
checks that perturbed cores were not copied from \(\beta_0\).  A decreasing
error window is recorded by
\[
    W_i=\ind\{e(h_i)>e(h_{i+1})>e(h_{i+2})\},
    \qquad i=1,2.
    \tag{P22-FD4}
\]
The status rule is
\[
\text{status}
=
\begin{cases}
\text{branch identity failure},&
    \min_{i,\pm} I_\pm(h_i)=0,\\
\text{copied-core failure},&
    \min_{i,\pm} K_\pm(h_i)=0,\\
\text{decreasing-window agreement},&
    \max(W_1,W_2)=1,\\
\text{inconclusive: no decreasing window},&
    \text{otherwise}.
\end{cases}
\tag{P22-FD5}
\]
The status is a derivative-consistency diagnostic for the fixed branch.  It is
not a posterior-accuracy theorem, a rank-stability theorem, or an HMC
convergence result.
\[
    \begin{aligned}
    \mathcal N_{\rm FD}
    =
    \{&
    \text{not exact posterior accuracy},\\
    &\text{not rank-stability certification},\\
    &\text{not HMC convergence certification}\}.
    \end{aligned}
    \tag{P22-FD7}
\]

A healthy derivative check typically shows decreasing error over a range of
\(h\), followed by roundoff or fitting-noise limitations.  If the error does
not decrease on any reasonable range, the derivative path is not yet aligned
with the value path.
\[
    \exists i\in\{1,2\}: W_i=1
    \quad\Longleftrightarrow\quad
    \text{there is a decreasing finite-difference error window.}
    \tag{P22-FD8}
\]

The local P22 branch identity condition is mandatory:
\[
    B(\beta_0)=B(\beta_0+h)=B(\beta_0-h)=B.
    \tag{P22-FD6}
\]
If the implementation rebuilds ranks, points, shifts, or domains at the
perturbed values, then the centered difference is no longer assessing the
derivative of \(\widehatell_2(\beta;B)\).  It is comparing a different
adaptive calculation.
If the implementation copies the numerical cores from \(\beta_0\) without
resolving the fixed fitting equations at the perturbed parameter, then it is
also assessing the wrong object.  The correct diagnostic freezes the branch choices but
allows all differentiable outputs of those choices to change with \(\beta\).
\[
    \begin{array}{ll}
    \min_{i,\pm}I_\pm(h_i)=0
    &\Rightarrow\ \text{the compared values use different branches},\\[2mm]
    \min_{i,\pm}K_\pm(h_i)=0
    &\Rightarrow\ \text{the perturbed values copied old numerical cores},\\[2mm]
    \max(W_1,W_2)=0
    &\Rightarrow\ \text{the diagnostic is inconclusive for derivative agreement}.
    \end{array}
    \tag{P22-FD9}
\]

\section{Minimal Mathematical Example}

The smallest useful mathematical example should be nonlinear but one-dimensional:
\[
    x_t=\beta x_{t-1}+\sigma\epsilon_t,
    \qquad
    \epsilon_t\sim N(0,1),
    \tag{P19-102}
\]
\[
    y_t=\sin(x_t)+\eta_t,
    \qquad
    \eta_t\sim N(0,r^2).
    \tag{P19-103}
\]
Use \(T=2\).  The first time step defines
\[
    q_1(x_1,x_0;\beta)
    =
    p_\beta(x_0)
    f_\beta(x_1\mid x_0)
    g_\beta(y_1\mid x_1).
    \tag{P19-104}
\]
It fits \(\phi_1\), computes
\[
    \widehat Z_1(\beta)=e^{-c_1}\int \phi_1(z_1,z_0;\beta)^2
    \dd z_1\dd z_0+\tau_1,
    \tag{P19-105}
\]
and stores
\[
    \widehat p_1(z_1;\beta)
    =
    \frac{\int
    \left[e^{-c_1}\phi_1(z_1,z_0;\beta)^2+\tau_1\lambda_1(z_1,z_0)\right]
    \dd z_0}
    {\widehat Z_1(\beta)}.
    \tag{P19-106}
\]
The second time step depends on that stored object:
\[
    q_2(x_2,x_1;\beta)
    =
    \widehat p_1(x_1;\beta)
    f_\beta(x_2\mid x_1)
    g_\beta(y_2\mid x_2).
    \tag{P19-107}
\]
The scalar checked by the example is
\[
    \widehatell_2(\beta)
    =
    \log\widehat Z_1(\beta)
    +
    \log\widehat Z_2(\beta).
    \tag{P19-108}
\]
The mathematical report object must contain
\[
    \widehatell_2(\beta_0),
    \qquad
    \partial_\beta\widehatell_2(\beta_0),
    \qquad
    D(h),
    \qquad
    |D(h)-\partial_\beta\widehatell_2(\beta_0)|.
    \tag{P19-109}
\]
It must also contain the branch manifest identifier used for \(\beta_0\),
\(\beta_0+h\), and \(\beta_0-h\).  The mathematical pass condition is that
those branch identifiers are identical and the derivative error has a stable
decreasing window before numerical roundoff dominates.

\section{Integrated Conclusion}

The fixed-branch gradient is built from five elementary ideas:
\[
\begin{array}{ll}
\text{normalizer derivative} & \partial\log Z=\dot Z/Z,\\[2mm]
\text{squared density derivative} & \partial(\phi^2)=2\phi\dot\phi,\\[2mm]
\text{mass contraction derivative} & \text{differentiate every factor once},\\[2mm]
\text{linear solve derivative} & N\dot g=\dot d-\dot N g,\\[2mm]
\text{carried filter derivative} & \partial(a/Z)=\dot a/Z-a\dot Z/Z^2.
\end{array}
\tag{P19-110}
\]
The tensor-train notation makes these ideas look large, but it does not change
their nature.  Once the branch is fixed, the forward pass defines one scalar.
The derivative pass differentiates that scalar.  What remains for empirical
work is to assess whether the approximation is accurate enough, stable enough,
and useful enough for the target scientific model.

\section{Reader Orientation Blocks For The Fixed-Branch Derivation}

This section collects neutral orientation blocks for the fixed-branch part.
They are not review gates.  They simply state, after the derivation, what has
been established and what a later implementation must preserve.

\subsection*{Squared-density block}

\[
\begin{array}{ll}
\text{established} &
    \widehat q_t=e^{-c_t}\phi_t^2+\tau_t\lambda_t\ge0,\\[1mm]
\text{stored} &
    \phi_t \text{ through } C_{t,k}
    \text{ and } \dot\phi_t \text{ through } \dot C_{t,k},\\[1mm]
\text{integrated} &
    e^{-c_t}\phi_t(z)^2+\tau_t\lambda_t(z),\\[1mm]
\text{differentiated} &
    e^{-c_t}\phi_t(z)^2
    \mapsto 2e^{-c_t}\phi_t(z)\dot\phi_t(z),\\[1mm]
\text{fixed} &
    c_t,\tau_t,\lambda_t,\text{ basis, domain, ranks, points},\\[1mm]
\text{failure mode} &
    \text{differentiating a different fitted square root than the one squared.}
\end{array}
\tag{P22-R1}
\]

\subsection*{Mass-contraction block}

\[
\begin{array}{ll}
\text{established} &
    R_t=\int\phi_t(z)^2\dd z
    \text{ is computed by TT mass contractions},\\[1mm]
\text{stored} &
    C_{t,k},\dot C_{t,k},B_{t,k},M_{>k},\dot M_{>k},\\[1mm]
\text{integrated} &
    \phi_t^2 \text{ one coordinate at a time},\\[1mm]
\text{differentiated} &
    \text{each product factor once, then summed},\\[1mm]
\text{fixed} &
    B_{t,k},\text{ coordinate order, rank pattern, basis family},\\[1mm]
\text{failure mode} &
    \text{omitting a product-rule term in } \dot M_{>k}.
\end{array}
\tag{P22-R2}
\]

\subsection*{Fixed-solve block}

\[
\begin{array}{ll}
\text{established} &
    N_k g_k=d_k
    \mapsto
    N_k\dot g_k=\dot d_k-\dot N_k g_k,\\[1mm]
\text{stored} &
    A_k,N_k,d_k,g_k \text{ and dotted counterparts},\\[1mm]
\text{integrated} &
    \text{nothing; this is a finite-dimensional fitting solve},\\[1mm]
\text{differentiated} &
    \text{the fixed normal equations},\\[1mm]
\text{fixed} &
    W,\rho,Z_{\rm fit},B_{\rm val},\text{sweep order},\\[1mm]
\text{failure mode} &
    \text{changing the solve system or copying }g_k(\beta_0)
    \text{ at perturbed parameters.}
\end{array}
\tag{P22-R3}
\]

\subsection*{Carried-filter block}

\[
\begin{array}{ll}
\text{established} &
    \widehat p_t=a_t/\widehat Z_t
    \text{ and }
    \dot{\widehat p}_t=\dot a_t/\widehat Z_t
    -a_t\dot{\widehat Z}_t/\widehat Z_t^2,\\[1mm]
\text{stored} &
    Q_t,\dot Q_t,P_t,\dot P_t
    \text{ with shape }(p,p),\\[1mm]
\text{integrated} &
    \widehat q_t(z_t,z_{t-1}) \text{ over } z_{t-1},\\[1mm]
\text{differentiated} &
    Q_t/\widehat Z_t \text{ by the quotient rule},\\[1mm]
\text{fixed} &
    \text{basis, query coordinate, branch, defensive term},\\[1mm]
\text{failure mode} &
    \text{feeding time }t+1\text{ without }\dot P_t
    \text{ or with an undeclared query rule.}
\end{array}
\tag{P22-R4}
\]

\appendix
\section{Source Coverage Summary}

The main text reconstructs the displayed mathematical spine of Zhao--Cui
Sections 1--3 and 5.  Equations (1)--(26) and (30)--(35), Algorithms 1--5,
the notation formulas, the square-root error inequalities, the upper and lower
conditional map formulas, the particle and smoothing weight normalizations, and
the preconditioning composition formulas are represented in BayesFilter
notation.  Section 4 error propagation and Section 6 numerical examples are
not reconstructed as the central object of this note; they are used only to
explain what the method does and does not prove.

\end{document}
